% Generated by report_gen.py (hwde v1) - do not hand-edit.
\documentclass[11pt]{article}
\newcommand\hspaper{a4paper}
\usepackage[nodiagramkit]{housestyle}
\usepackage{longtable}
\usepackage{pdfpages}
\makeatletter
\@ifundefined{KV@Gin@artifact}{\define@key{Gin}{artifact}[]{}}{}
\makeatother
\setcounter{tocdepth}{1}
\hsslug{PCB-0017-A \textperiodcentered\ stereo-class-d-amp}
\renewcommand{\sectionmark}[1]{\markboth{#1}{}}
\hssection{\leftmark}
\begin{document}
\sloppy
\hstitleblock{PCB-0017-A \textperiodcentered\ hwde v1 pipeline \textperiodcentered\ generated 2026-09-29 13:03}%
{stereo-class-d-amp \textemdash{} Full Design Document}%
{Phase P7; 4/5 gates pass (not passing: verify).}
\tableofcontents

\section{Board and Run Metadata}
{\small
\begin{longtable}{lp{11cm}}
\hstoprule
\hshead{Field} & \hshead{Value} \\
\hline
\endhead
board & stereo-class-d-amp \\
part number & PCB-0017-A \\
workspace & \texttt{boards/PCB-0017-A\_stereo-class-d-amp} \\
phase & P7 \\
gate status & 4/5 gates pass (not passing: verify) \\
state created & 2026-09-27T00:25:28 \\
state updated & 2026-09-27T08:57:51 \\
\hline
\end{longtable}
}
\section{Overview}
\subsection*{PCB-0017-A stereo-class-d-amp ? brief}
Owner brief: rebuild and compare against a published 44 x 34 mm reference design.

\begin{itemize}
\item TPA3116D2 stereo class-D amplifier on a single 12 V supply, stereo input and output, SMD parts only,
2-layer board as small as possible. Target: match or beat the reference design's 44 x 34 mm.
\item Correctness and the TPA3116D2's thermal pad / heat-sinking come before size; say where we deviate.
\item The reference design's reported parts: TPA3116D2; 2.2 uF X7R input coupling; 100 uF + 100 nF + 1 nF supply bypass;
10 uH LC output filters per channel with damping branches; 5.6k/100k config and pull-up resistors;
wire pads J1-J4 and TP1. About 40 SMD parts, 45 footprints.
\item Source: https://jlcpcb.com/blog/gpt-6-astra-pcb-design-in-kicad
\item Stop at a DRC/ERC-clean fab package. Do NOT order. End with COMPARISON.md: our size, layers,
part count and DRC result against the reference design's.
\end{itemize}
\section{Requirements}
\subsection*{PCB-0017-A stereo-class-d-amp - requirements}
Source: \texttt{brief/brief.md} (owner brief). No mode token: normal design, full scope. Run is unattended; every open question in section 9 carries a default the orchestrator may take as the owner's delegate.

\subsubsection*{1. Function}
A stereo class-D audio amplifier board around the TI TPA3116D2, fed from a single 12 V supply, with stereo line-level input and stereo speaker output on wire pads. It reproduces a published 44 x 34 mm reference design (JLCPCB blog post, link in the brief) so we can compare our result against the reference design's: a 2-layer board, SMD parts only, as small as possible, with correctness and the TPA3116D2's thermal pad and heat-sinking taking priority over size. Work stops at a DRC/ERC-clean fab package (no order) plus a \texttt{COMPARISON.md} giving our size, layer count, part count and DRC result against the reference design's reported 44 x 34 mm, 2 layers, about 40 SMD parts and 45 footprints. Where we deviate from the reference design we say where and why.

\subsubsection*{2. Interfaces}
\begin{itemize}
\item Audio input, stereo, line level, single-ended: two channels plus ground, AC-coupled into the
TPA3116D2 through 2.2 uF X7R capacitors (the reference design's value). Connector: wire pads (the reference design's J1-J4
are wire pads, not connectors). Input impedance and level are the amplifier's defaults
(gain set by the 5.6k / 100k resistors the reference design reports). Which of J1-J4 is input vs output vs
power is not stated - see section 9 Q1.
\item Audio output, stereo, bridge-tied load (BTL) speaker outputs, one 10 uH LC output filter per
channel with a damping branch (the reference design's parts). Wire pads. Speaker impedance not stated - Q2.
\item Power input, 12 V DC, wire pads. Q3 covers a screw terminal or barrel jack instead.
\item TP1: one test point (the reference design's list). Purpose not stated; ASSUMED: a ground or supply probe pad.
\item Configuration: gain / mode set by fixed resistors (the reference design's 5.6k / 100k). No switches, no LEDs,
no buttons stated. Fault / shutdown / mute pins: pull-up resistors only (the reference design's list), no user
control - Q4 asks whether a mute or fault indicator is wanted.
\item No USB, no RF, no digital interface.
\end{itemize}
\subsubsection*{3. Power}
\begin{itemize}
\item Input: single 12 V DC supply (brief). Range not stated; ASSUMED 12 V nominal, tolerate roughly
10.8 to 13.2 V (a regulated 12 V adapter). The TPA3116D2 accepts a wide range, but the input
connector, bulk capacitor voltage rating and reverse-polarity choice hinge on the range - Q5.
\item Rails: PVCC = 12 V direct to the amplifier. The TPA3116D2 generates its own internal analog and
gate-drive rails; no external regulator is expected. GUESS: no other rail on the board.
\item Supply bypass per the reference design: 100 uF bulk + 100 nF + 1 nF at each PVCC group (the reference design reports one set;
the datasheet-driven count per pin group is P1's to settle).
\item Supply current GUESS (not researched, from the part's class): at 12 V the TPA3116D2 delivers on
the order of 15 W per channel into 4 ohm or about 8 W per channel into 8 ohm at the onset of
clipping. Both channels driven at full output into 4 ohm draws about 2.5 to 3 A from 12 V;
typical music program is well under 1 A. Whether this is treated as \textgreater{}3 A is a section 8 flag - Q6.
\item No battery, no charging.
\item Reverse-polarity or overvoltage protection at the input: not stated - Q7.
\end{itemize}
\subsubsection*{4. Environment}
Not stated. ASSUMED: indoor bench / hobby use, 0 to 40 C ambient, no enclosure specified, no ingress rating, no vibration requirement. The TPA3116D2 dissipates a few watts at full output, so thermal design assumes free air and relies on the PCB copper (thermal pad vias to a bottom copper pour) rather than an attached heatsink - Q8 confirms no clip-on heatsink is required.

\subsubsection*{5. Size \& mounting}
\begin{itemize}
\item Outline: as small as possible; target to match or beat the reference design's 44 x 34 mm on 2 layers. This is a
soft target, no HARD cap: correctness and the thermal pad / heat-sinking area come first, and the
brief asks us to state where we deviate. If the thermal copper needs more area than 44 x 34 mm
allows, the board grows and COMPARISON.md says so.
\item Layers: 2 (brief). Treated as a firm design choice, not a physical cap; P1 may argue for more only
with a stated reason and the deviation recorded.
\item Mounting holes: not stated. The reference design's board is not reported to have any. ASSUMED none (bare module
with wire pads), see Q9.
\item Height limit: none stated. The 100 uF bulk capacitor is the tallest part; SMD only per the brief.
\item SMD parts only (brief). Wire pads J1-J4 and TP1 are copper pads, no through-hole connectors.
\end{itemize}
\subsubsection*{6. Quantity \& budget}
Not stated. Nothing is ordered in this run (brief: stop at the fab package). ASSUMED for costing and DFM: 5 boards, hobby budget, no unit-cost target - Q10.

\subsubsection*{7. Assembly}
Not stated. ASSUMED: JLC PCBA, all parts on the top side (single-sided assembly) where the layout allows; the bottom side kept for the thermal copper pour and vias only. A double-sided assembly is allowed if it is the only way to hit the size target, and COMPARISON.md then says so - Q11. Parts should be JLC-stocked (basic / preferred where a choice exists) so the package is orderable as-is even though we do not order.

\subsubsection*{8. Compliance/safety flags}
\begin{itemize}
\item Mains voltage: no. The board takes 12 V DC from an external supply.
\item Batteries: no.
\item Motors: no.
\item \textgreater{}30 V: no (12 V supply; the TPA3116D2's own switching stays under 30 V).
\item High current (\textgreater{}3 A): borderline - flagged. The 12 V input and the BTL speaker outputs carry
amplifier current. At 12 V into 4 ohm, both channels at clipping, the supply draw is estimated
near 3 A (section 3 GUESS); the output filter inductors and traces carry a few amps peak per
channel. Q6 asks how to treat it. Not a guess to proceed on: the trace widths, inductor
saturation rating and pad sizing depend on the answer.
\item RF transmit: no. Class-D switching (about 400 kHz to 1.2 MHz) radiates; the LC output filter
and layout handle EMI but no intentional transmitter is present.
\item Speaker output is a bridge-tied load: neither output terminal is ground. Note for the wire pad
labelling and for anyone wiring it; not a safety hazard at these voltages.
\end{itemize}
\subsubsection*{9. Open questions}
1. Which of the wire pads J1-J4 carry what? Default: J1 = 12 V power in (V+ and GND), J2 = stereo line input (L, R, GND), J3 = left speaker out (+/-), J4 = right speaker out (+/-). Each "pad" is a small group of copper pads. 2. Speaker impedance the amplifier should be designed for? Default: 4 ohm minimum, 8 ohm typical (the LC output filter values are tuned for this; 10 uH per the reference design suits both). 3. Should the power input stay as bare wire pads like the reference design's, or be a real connector (screw terminal or barrel jack)? Default: wire pads, to match the reference design and the SMD-only rule. 4. Do you want a mute/standby control or a fault indicator LED? Default: no - fixed pull-ups so the amp is always on, matching the reference design's parts list. 5. Supply voltage range the board must accept? Default: 12 V nominal, 10.8 to 13.2 V (a regulated 12 V adapter); capacitors rated 25 V or better. 6. Current handling (safety flag, section 8): design the 12 V input path and speaker outputs for how much continuous current? Default: 3 A continuous on the 12 V input, 3 A per speaker output, which covers both channels driven fully into 4 ohm at 12 V. Answering "yes to the default" accepts that the board is sized for up to 3 A; a higher figure widens traces and grows the board. 7. Reverse-polarity protection on the 12 V input? Default: none (matches the reference design's about-40-part count); the user must wire polarity correctly. 8. Is PCB copper alone acceptable for heat-sinking (thermal pad vias to a bottom copper pour), with no attached heatsink? Default: yes - the TPA3116D2 at 12 V into 8 ohm runs cool on copper alone; into 4 ohm continuously it may need more copper area than 44 x 34 mm, and we would then grow the board and record the deviation. 9. Mounting holes? Default: none (bare module, like the reference design's). 10. Quantity and budget for the DFM / cost sanity check? Default: 5 boards, no unit-cost target, JLC-stocked parts. 11. Assembly: JLC PCBA, top side only? Default: yes, single-sided top; allow bottom-side parts only if that is the sole way to meet the size target, and say so in COMPARISON.md. 12. Input level and gain: line level (about 1 Vrms) with the gain the reference design's 5.6k/100k resistors set? Default: yes, keep the reference design's gain setting so the comparison is like-for-like.

\paragraph*{Answers to section 9 (orchestrator as owner's delegate, 2026-09-27; recorded as state.json decisions)}
All twelve defaults above are ACCEPTED as written. Q6 (current) is PROVISIONAL: 3 A continuous on the 12 V input and 3 A per speaker output; re-confirmed at H1 against the P1 power research.

\section{Architecture}
\subsection*{Decisions of record (state.json)}
{\small
\begin{longtable}{lp{5.2cm}p{7.2cm}}
\hstoprule
\hshead{Phase} & \hshead{What} & \hshead{Why} \\
\hline
\endhead
P0 & Q1 pads: J1=12V in (V+,GND), J2=stereo line in (L,R,GND), J3=left spk +/-, J4=right spk +/-; TP1 per the reference design & delegate answer (owner said decide yourself); matches the reference design's J1-J4 \\
P0 & Q2 speaker: 4 ohm min, 8 ohm typ & delegate; 10uH filter suits both \\
P0 & Q3 power in: bare wire pads & delegate; like-for-like with the reference design \\
P0 & Q4 no mute/fault LED; fixed pull-ups, always on & delegate; the reference design has none \\
P0 & Q5 supply 12V nom, 10.8-13.2V, caps \textgreater{}=25V & delegate default \\
P0 & Q6 PROVISIONAL (safety, \textgreater{}3A borderline): 3A continuous on 12V input, 3A per speaker output & delegate; TPA3116D2 \textasciitilde{}2x15W into 4ohm at 12V -\textgreater{} \textasciitilde{}2.8A supply; unattended run, reconfirm at H1 \\
P0 & Q7 no reverse-polarity protection & delegate; matches the reference design part count; noted as a deviation risk in COMPARISON \\
P0 & Q8 PCB copper heat-sinking only (PowerPAD thermal vias to bottom pour), grow board if needed and record & owner: thermal first, then size \\
P0 & Q9 no mounting holes & delegate; bare module like the reference design \\
P0 & Q10 qty 5, no unit-cost target, JLC stocked parts & delegate default \\
P0 & Q11 JLC PCBA top side only unless size forces otherwise & delegate default \\
P0 & Q12 line-level in, keep the reference design's 5.6k/100k gain/config & delegate; like-for-like \\
P1 & DEVIATION: amp IC TPA3116D2DAD -\textgreater{} TPA3118D2DAP (same SLOS708 family, same 32-pin HTSSOP pinout, PowerPAD DOWN) & TPA3116D2 exists only as DAD (pad UP, TI SLOS708); its 14 C/W needs a top heatsink, which the reference design's own write-up admits is not modeled. Brief: thermal pad/heat-sinking first, SMD only. At 12 V output is rail-limited so no audio power lost; pad-down lets 2-layer copper + vias sink it (\textasciitilde{}22-30 C/W). \\
P1 & Q6 revised: +12V input and its GND return sized for 4 A; 3 A per speaker output; inductor Isat \textgreater{}= 4 A & power research: 3.6 A at 10\% THD into 4 ohm, 4.0 A at 13.2 V \\
P1 & Sustained full-power 4 ohm sine at 40 C is NOT a requirement (music + 8 ohm full power are); keep the size target & power research: that case needs \textasciitilde{}50x50 mm 2 oz or PLIMIT; thermal shutdown is the chip's own protection; FAULTZ-\textgreater{}SDZ gives auto-recovery \\
P1 & Prefer 2 oz outer copper if the stackup table offers it for 2L & halves trace widths and improves spreading; 0.65 mm pitch is fine at 2 oz floors \\
P1 & PLIMIT tied to GVDD (no limit); FAULTZ tied to SDZ with 100k pull-up for self-clearing thermal fault; gain 20 dB via 5.6k master & TI 12 V design; refdesign research resolves the reference design's 5.6k/100k as gain leg + pull-ups \\
P1 & No reverse-polarity/TVS part (Q7 stands) & like-for-like with the reference design; risk (reversed supply, hot-plug ring) written in COMPARISON.md \\
P1 & 10 uF PVCC decoupling in 0805 (Basic CL21A106KAYNNNE), not 1206 & Basic part, smaller footprint helps size; 25 V rating kept \\
P1 & No switch-node EMI snubber (10R+330pF); keep only the 3.3R+10nF damping branch per channel output & like-for-like with the reference design and the EVM damping network; EMI not a stated requirement \\
P1 & Parts leads: TPA3118D2DAPR C46497; SLO0630H100MTT 10uH C207842 (Isat 5.5A); 680nF out filter; bulk MA25V100M6x6 polymer + RVT 100uF 35V electrolytic & P1 component-scout shortlist \\
P2 & P2 research tiers: class-d-amp + lc-output-filter researchers on fable/high (tier); dc-input + audio-line-in researchers and second readers on sonnet/high (explicit downgrade) & B1/B2 are low-risk textbook blocks already covered by P1 power.md/refdesign; owner asked for quality without burning money \\
P2 & Start P3 part-sourcer in parallel with P2 research, before H1 & part choice doesn't turn on the checklist research; a value the research changes is amended in parts.json before P4 \\
P3 & C11/C12 10 uF PVCC: accept X5R CL21A106KAYNNNE (C15850) though sheets.md said X7R & caps sit below 85 C (X5R limit) at \textless{}=1.5 W design dissipation; DC-bias loss at 12 V is covered by the 100 uF bulk per pair; Basic part \\
P3 & R4-R7 3.3R 0603 100 mW (RC0603FR-073R3L) accepted & branch dissipation \textasciitilde{}1 mW: 10 nF is \textasciitilde{}800 ohm at 20 kHz and \textasciitilde{}40 ohm at 400 kHz against \textless{}0.3 V residual ripple \\
P2 & R4-R7+C23-C26 (3.3R+10nF) stay on the speaker side of the LC filter as TI's EMI C-RC snubber, not at the switch node & research: TI labels it an EMI snubber (corner \textasciitilde{}4.8 MHz), not audio damping; at the switch node 10 nF x 12 V\textasciicircum{}2 x 400 kHz = 0.58 W would overload a 0603. Filter Q 0.52-1.04 over 4-8 ohm is acceptable per SLAA701B \\
P3 & designing under coverage gap: block:B1 dc-input classes power-loop, esd, sourcing, constraints-emission have no record at the approved floor; no second research round & B1 is two wire pads + two 100uF bulk caps with no protection by decision Q7; its loop/placement rules come from the verified B3 PVCC three-tier decoupling and dc-input hot-plug records; a 2nd round would cost \textasciitilde{}\$5-8 for little design delta \\
P3 & no TI Fig 37 PVCC RC damper (C58 100nF + R10 3.3R) added & C14 electrolytic ESR is the bulk damping per the verified dc-input-bulk-esr-damping-budget record; hot-plug ring margin 1.15x stays a P4-review watch item with RC snubber as fallback \\
P5 & U1 footprint fix: the 9 EasyEDA thermal vias (0.3 drill, 0.61 pad, 1.0 pitch 3x3) had empty pad numbers (no net -\textgreater{} 9 mask-bridge + 9 clearance errors); renumbered to pad 33 (GND) and set layers {*}.Cu only (no paste/mask on the vias: bottom tented, top inside the EP mask window) & thermal vias must be GND and join the bottom pour; record class-d-amp-powerpad-thermal-via-array asks component-side capping, which is impossible inside the EP window without plugging - via-in-pad plugging is the P9 alternative; 3x3 at 1.0 mm meets the \textgreater{}=6 count / \textless{}=1.2 mm pitch rule \\
P5 & wire-pad footprints: moved pin-1 silk corner 0.35 mm outside pad 1 (was over copper, 4 silk\_over\_copper warnings) & footprint-internal silk is a librarian fix; board\_init self-check now 0 violations \\
P4 & review W1/W2: swap C13 to a \textgreater{}=35 V low-ESR polymer/hybrid D6.3 part and C14 to a 35 V part with a stated ESR/ripple; ADD a unidirectional TVS (SMBJ15A-class, clamp \textless{} 30 V abs max) across J1 +12V/GND & reviewer: 26 V hot-plug ring exceeds C13's 25 V and C14's ESR cannot damp the bank; TVS bounds the ring below U1's 30 V abs max. Correctness over the reference design's part count (deviation for COMPARISON) \\
P4 & review W3: C1-C4 2.2 uF -\textgreater{} 1 uF X7R 50 V 0805, K tolerance (TI Fig 37 value; corner 2.7 Hz at 60k) & 2.2 uF x 60k (132 ms) with no source connected reaches the 420 ms DC-detect window at tolerance corners and fault-cycles on power-up; deviation from the reference design's 2.2 uF recorded \\
P4 & review W4: reverse polarity stays unprotected (Q7) but J1 gets a +12V/GND silk legend; W5: C11/C12 X5R 10 uF accepted & W4 accepted risk, legend is free; W5: \textasciitilde{}4-5 uF effective at 12 V is ample beside 200 uF bulk + 100 nF/1 nF HF caps \\
P4 & Schematic review W1-W3 fixed: C13 35V polymer C46550467, C14 FZ 35V C2992596 (stated ESR/ripple), C1-C4 1uF X7R C28323, D1 SMBJ15A C113988 across J1 & W1 26V ring \textgreater{} 25V C13 and C14 cannot damp it -\textgreater{} TVS clamps; W2 RVT had no ESR spec; W3 132ms input tau vs 420ms DC-detect -\textgreater{} 60ms \\
P6 & Accept 53.6 x 36.8 mm outline (2 layers) as a deviation from the 44 x 34 target & Thermal/decoupling layout first: 25x25 mm clear thermal copper around U1, PVCC stacks at the pins, 7.3 mm inductors need 9.6 mm column pitch and 3x6 mm speaker wire pads J3/J4 take their own column; closing \textasciitilde{}9 mm would need smaller inductors or wire pads, trading current rating/solderability for size \\
P6 & Strip body-outline silk from C0402/C0603/C0805/R0603 (lib + board) and widen IND courtyard to cover pads & 98 outline-vs-outline silk\_overlap warnings in the dense PVCC/bootstrap stacks would fail drc\_routed; passives are unpolarised, refdes kept; IND courtyard was smaller than its pads \\
P7 & P7-\textgreater{}P6: spread L3/L4 up and L1/L2 down \textasciitilde{}2.5 mm (grow outline height if needed) to fit the C15-C18 bootstrap caps in two columns and route OUT nets on F.Cu & OUT nets unroutable with current placement; caps-on-B.Cu alternative would cut the U1 heat-spreader pour edge with 8 vias; thermal first, area already a stated deviation \\
\hline
\end{longtable}
}
\subsection*{Sheet plan}
\subsection*{PCB-0017-A stereo-class-d-amp - schematic sheet plan (P2)}
\subsubsection*{Ruling: ONE sheet, no hierarchy}
38 fitted parts around one IC, on one signal path (J2 -\textgreater{} U1 -\textgreater{} LC -\textgreater{} J3/J4) with one rail. A hierarchy would buy nothing but hierarchical pins and sheet-path net names; a single A3-class sheet shows the whole amplifier at a glance and the placement groups at P6 read straight off it. Root sheet = the only sheet: \texttt{PCB-0017-A\_stereo-class-d-amp.kicad\_sch} (same name as the KiCad project and the workspace directory, boards CLAUDE.md).

Net naming: power nets bare (\texttt{+12V}, \texttt{GND}); every other net is a local label and carries a leading \texttt{/} in the netlist (\texttt{/OUTPL}). These are the CANONICAL names: \texttt{constraints.json} uses them and the schematic must produce them exactly.

\subsubsection*{Drawing layout (left to right, mirroring the board)}
J2 and the four input caps at the far left; U1 in the middle drawn with pins 1-16 on its left and 17-32 on its right (the DAP footprint's own orientation: pin 1 top-left, PVCC 31/32 top-right, AVCC/PVCC 17-19 bottom-right); the configuration cluster (R1, R2, R3, C5, the MODSEL/AM/SYNC ties) below-left of U1; bootstrap and PVCC decoupling caps drawn at their pins on the right; the two LC filters to the right of U1 with the right channel above the left (matching pins 27/29 above 21/23); J4 then J3 at the far right; J1 with C13/C14 along the top right; TP1 on GND anywhere; one PWR\_FLAG each on \texttt{+12V} and \texttt{GND} (external rail, no on-board source - ERC needs them).

\subsubsection*{Pin-by-pin tie table (TI SLOS708G, DAP package - P3/P4 verify against the symbol)}
\begin{quote}\small\ttfamily\raggedright
\textbar{} pin \textbar{} name \textbar{} net \textbar{} tie / part \textbar{}\\
\textbar{}---\textbar{}---\textbar{}---\textbar{}---\textbar{}\\
\textbar{} 1 \textbar{} MODSEL \textbar{} `GND` \textbar{} BD mode \textbar{}\\
\textbar{} 2 \textbar{} SDZ \textbar{} `/SDZ` \textbar{} joined with FAULTZ; R2 100k to `+12V` \textbar{}\\
\textbar{} 3 \textbar{} FAULTZ \textbar{} `/SDZ` \textbar{} open-drain, same node as SDZ (auto-clearing fault) \textbar{}\\
\textbar{} 4 \textbar{} RINP \textbar{} `/RINP` \textbar{} C3 2.2 uF from `/IN\_R` \textbar{}\\
\textbar{} 5 \textbar{} RINN \textbar{} `/RINN` \textbar{} C4 2.2 uF to `GND` \textbar{}\\
\textbar{} 6 \textbar{} PLIMIT \textbar{} `/GVDD` \textbar{} direct tie, no limit \textbar{}\\
\textbar{} 7 \textbar{} GVDD \textbar{} `/GVDD` \textbar{} C5 1 uF to `GND` \textbar{}\\
\textbar{} 8 \textbar{} GAIN/SLV \textbar{} `/GAIN` \textbar{} R1 5.6k to `GND`, GVDD leg open: 20 dB master \textbar{}\\
\textbar{} 9 \textbar{} GND \textbar{} `GND` \textbar{} \textbar{}\\
\textbar{} 10 \textbar{} LINP \textbar{} `/LINP` \textbar{} C1 2.2 uF from `/IN\_L` \textbar{}\\
\textbar{} 11 \textbar{} LINN \textbar{} `/LINN` \textbar{} C2 2.2 uF to `GND` \textbar{}\\
\textbar{} 12 \textbar{} MUTE \textbar{} `/MUTE` \textbar{} R3 100k to `GND` (outputs enabled) \textbar{}\\
\textbar{} 13 \textbar{} AM2 \textbar{} `GND` \textbar{} 400 kHz (AM2/AM1/AM0 = 0/0/0) \textbar{}\\
\textbar{} 14 \textbar{} AM1 \textbar{} `GND` \textbar{} \textbar{}\\
\textbar{} 15 \textbar{} AM0 \textbar{} `GND` \textbar{} \textbar{}\\
\textbar{} 16 \textbar{} SYNC \textbar{} no connect \textbar{} output in master mode; ERC no-connect flag \textbar{}\\
\textbar{} 17 \textbar{} AVCC \textbar{} `+12V` \textbar{} C6 100 nF to `GND` at the pin \textbar{}\\
\textbar{} 18, 19 \textbar{} PVCC \textbar{} `+12V` \textbar{} C7 1 nF, C9 100 nF, C11 10 uF, C14 100 uF electrolytic \textbar{}\\
\textbar{} 20 \textbar{} BSNL \textbar{} `/BSNL` \textbar{} C15 220 nF to `/OUTNL` \textbar{}\\
\textbar{} 21 \textbar{} OUTNL \textbar{} `/OUTNL` \textbar{} L2 -\textgreater{} `/SPK\_L\_N` \textbar{}\\
\textbar{} 22 \textbar{} GND \textbar{} `GND` \textbar{} \textbar{}\\
\textbar{} 23 \textbar{} OUTPL \textbar{} `/OUTPL` \textbar{} L1 -\textgreater{} `/SPK\_L\_P` \textbar{}\\
\textbar{} 24 \textbar{} BSPL \textbar{} `/BSPL` \textbar{} C16 220 nF to `/OUTPL` \textbar{}\\
\textbar{} 25 \textbar{} GND \textbar{} `GND` \textbar{} \textbar{}\\
\textbar{} 26 \textbar{} BSNR \textbar{} `/BSNR` \textbar{} C17 220 nF to `/OUTNR` \textbar{}\\
\textbar{} 27 \textbar{} OUTNR \textbar{} `/OUTNR` \textbar{} L4 -\textgreater{} `/SPK\_R\_N` \textbar{}\\
\textbar{} 28 \textbar{} GND \textbar{} `GND` \textbar{} \textbar{}\\
\textbar{} 29 \textbar{} OUTPR \textbar{} `/OUTPR` \textbar{} L3 -\textgreater{} `/SPK\_R\_P` \textbar{}\\
\textbar{} 30 \textbar{} BSPR \textbar{} `/BSPR` \textbar{} C18 220 nF to `/OUTPR` \textbar{}\\
\textbar{} 31, 32 \textbar{} PVCC \textbar{} `+12V` \textbar{} C8 1 nF, C10 100 nF, C12 10 uF, C13 100 uF polymer \textbar{}\\
\textbar{} 33 \textbar{} PowerPAD \textbar{} `GND` \textbar{} via array to B.Cu \textbar{}
\end{quote}
\subsubsection*{Nets that leave a block (placement groups at P6, net classes at P5)}
\begin{quote}\small\ttfamily\raggedright
\textbar{} net \textbar{} from \textbar{} to \textbar{} class / note \textbar{}\\
\textbar{}---\textbar{}---\textbar{}---\textbar{}---\textbar{}\\
\textbar{} `+12V` \textbar{} J1.1 \textbar{} C13, C14, C11, C12, C9, C10, C7, C8, U1 17/18/19/31/32, R2 \textbar{} 4 A, \textgreater{}= 1.0 mm at 2 oz \textbar{}\\
\textbar{} `GND` \textbar{} J1.2, J2.2, TP1 \textbar{} everything; B.Cu pour + F.Cu pour \textbar{} 4 A return, heat spreader \textbar{}\\
\textbar{} `/IN\_L`, `/IN\_R` \textbar{} J2.3, J2.1 \textbar{} C1, C3 \textbar{} signal, keep \textgreater{}= 10 mm from L1-L4 \textbar{}\\
\textbar{} `/LINP`, `/RINP` \textbar{} C1, C3 \textbar{} U1 10, 4 \textbar{} 60 kOhm high-Z nodes, short \textbar{}\\
\textbar{} `/LINN`, `/RINN` \textbar{} C2, C4 \textbar{} U1 11, 5 \textbar{} high-Z, short \textbar{}\\
\textbar{} `/GVDD` \textbar{} U1 7 \textbar{} U1 6, C5 \textbar{} local \textbar{}\\
\textbar{} `/GAIN` \textbar{} U1 8 \textbar{} R1 \textbar{} local; sets gain and master mode, latched at start-up \textbar{}\\
\textbar{} `/SDZ` \textbar{} U1 2, 3 \textbar{} R2 \textbar{} local \textbar{}\\
\textbar{} `/MUTE` \textbar{} U1 12 \textbar{} R3 \textbar{} local \textbar{}\\
\textbar{} `/OUTPL`, `/OUTNL` \textbar{} U1 23, 21 \textbar{} L1, L2, C16, C15 \textbar{} 3 A switch nodes, \textgreater{}= 0.8 mm, symmetric pair, tight \textbar{}\\
\textbar{} `/OUTPR`, `/OUTNR` \textbar{} U1 29, 27 \textbar{} L3, L4, C18, C17 \textbar{} same \textbar{}\\
\textbar{} `/BSPL`, `/BSNL`, `/BSPR`, `/BSNR` \textbar{} U1 24, 20, 30, 26 \textbar{} C16, C15, C18, C17 \textbar{} bootstrap, loop local to its pin pair \textbar{}\\
\textbar{} `/SPK\_L\_P`, `/SPK\_L\_N` \textbar{} L1, L2 \textbar{} C19, C20, R4, R5, J3.1, J3.2 \textbar{} 3 A, \textgreater{}= 0.8 mm \textbar{}\\
\textbar{} `/SPK\_R\_P`, `/SPK\_R\_N` \textbar{} L3, L4 \textbar{} C21, C22, R6, R7, J4.1, J4.2 \textbar{} 3 A \textbar{}\\
\textbar{} `/ZL\_P`, `/ZL\_N`, `/ZR\_P`, `/ZR\_N` \textbar{} R4, R5, R6, R7 \textbar{} C23, C24, C25, C26 \textbar{} damping-branch midpoints, local \textbar{}
\end{quote}
\subsubsection*{Refdes ranges (one sheet; fixed so a later renumber lands in the same places)}
\begin{quote}\small\ttfamily\raggedright
\textbar{} class \textbar{} range \textbar{} assigned \textbar{}\\
\textbar{}---\textbar{}---\textbar{}---\textbar{}\\
\textbar{} U \textbar{} U1-U9 \textbar{} U1 TPA3118D2DAPR \textbar{}\\
\textbar{} L \textbar{} L1-L9 \textbar{} L1 OUTPL, L2 OUTNL, L3 OUTPR, L4 OUTNR - 10 uH 7.3 x 6.6 mm \textbar{}\\
\textbar{} R \textbar{} R1-R19 \textbar{} R1 5.6k gain; R2 100k SDZ pull-up; R3 100k MUTE pull-down; R4 SPK\_L\_P, R5 SPK\_L\_N, R6 SPK\_R\_P, R7 SPK\_R\_N damping 3.3R 0603 \textbar{}\\
\textbar{} C \textbar{} C1-C39 \textbar{} C1 LINP / C2 LINN / C3 RINP / C4 RINN 2.2 uF 0805; C5 1 uF GVDD 0603; C6 100 nF AVCC 0603; C7, C8 1 nF C0G 0603 (PVCC 18/19, 31/32); C9, C10 100 nF 0603; C11, C12 10 uF 0805; C13 100 uF polymer (31/32), C14 100 uF electrolytic 35 V (18/19); C15 BSNL, C16 BSPL, C17 BSNR, C18 BSPR 220 nF 0805; C19 SPK\_L\_P, C20 SPK\_L\_N, C21 SPK\_R\_P, C22 SPK\_R\_N 680 nF 0805; C23-C26 10 nF 0402 damping (same order) \textbar{}\\
\textbar{} J \textbar{} J1-J9 \textbar{} J1 power 2 pads; J2 line in 3 pads (R, GND, L); J3 left spk 2 pads; J4 right spk 2 pads - all SMD wire pads, no connector body \textbar{}\\
\textbar{} TP \textbar{} TP1-TP9 \textbar{} TP1 GND \textbar{}\\
\textbar{} `\#PWR` \textbar{} `pwr\_base = 100` -\textgreater{} `\#PWR0101...` \textbar{} `+12V`, `GND` symbols \textbar{}\\
\textbar{} `\#FLG` \textbar{} `\#FLG0101...` \textbar{} one PWR\_FLAG on `+12V`, one on `GND` \textbar{}
\end{quote}
\subsubsection*{Values and ratings the schematic must carry}
\begin{quote}\small\ttfamily\raggedright
\textbar{} ref \textbar{} value \textbar{} rating / dielectric \textbar{} why \textbar{}\\
\textbar{}---\textbar{}---\textbar{}---\textbar{}---\textbar{}\\
\textbar{} C1-C4 \textbar{} 2.2 uF \textbar{} X7R, 50 V (25 V min), 0805 \textbar{} the reference design's value; 1.2 Hz corner at 60 k \textbar{}\\
\textbar{} C5 \textbar{} 1 uF \textbar{} X7R 25 V 0603 \textbar{} GVDD LDO decoupling, datasheet-mandated \textbar{}\\
\textbar{} C6 \textbar{} 100 nF \textbar{} X7R 50 V 0603 \textbar{} AVCC \textbar{}\\
\textbar{} C7, C8 \textbar{} 1 nF \textbar{} C0G/NP0 50 V 0603 \textbar{} switching-edge loop \textbar{}\\
\textbar{} C9, C10 \textbar{} 100 nF \textbar{} X7R 50 V 0603 \textbar{} mid-frequency \textbar{}\\
\textbar{} C11, C12 \textbar{} 10 uF \textbar{} X7R 25 V 0805 \textbar{} P1 decision (Basic 0805, not 1206) \textbar{}\\
\textbar{} C13 \textbar{} 100 uF \textbar{} Al-polymer 25 V, ESR \textless{}= 30 mOhm, \textgreater{}= 2 A rms \textbar{} ripple \textbar{}\\
\textbar{} C14 \textbar{} 100 uF \textbar{} Al-electrolytic 35 V low-ESR \textbar{} ripple + damps hot-plug ring \textbar{}\\
\textbar{} C15-C18 \textbar{} 220 nF \textbar{} X7R 50 V 0805 \textbar{} bootstrap, one per output \textbar{}\\
\textbar{} C19-C22 \textbar{} 680 nF \textbar{} X7R 50 V 0805 \textbar{} LC filter, f0 61 kHz with 10 uH \textbar{}\\
\textbar{} C23-C26 \textbar{} 10 nF \textbar{} X7R 50 V 0402 \textbar{} damping branch \textbar{}\\
\textbar{} L1-L4 \textbar{} 10 uH \textbar{} Isat \textgreater{}= 4 A (5.5 A lead), Irms \textgreater{}= 3 A, DCR \textless{}= 70 mOhm, 7.3 x 6.6 mm shielded \textbar{} 3.5 A peak, 7.5 A OC trip \textbar{}\\
\textbar{} R1 \textbar{} 5.6k \textbar{} 1 \%, 0603 \textbar{} 20 dB / master; latched at power-up \textbar{}\\
\textbar{} R2, R3 \textbar{} 100k \textbar{} 1 \%, 0603 \textbar{} SDZ+FAULTZ pull-up, MUTE pull-down \textbar{}\\
\textbar{} R4-R7 \textbar{} 3.3R \textbar{} 1 \%, 0603 (\textgreater{}= 100 mW) \textbar{} damping \textbar{}
\end{quote}
Full architecture narrative (not inlined; see the artifact index): \texttt{architecture/blocks.md}, \texttt{architecture/power\_tree.md}, \texttt{architecture/stackup.md}.
\section{Schematic}
Gate \texttt{erc} (phase P4): \textbf{pass} after 2 attempt(s); last 2026-09-27T03:45:59, failing 0 of 0.
The full schematic PDF follows.
\includepdf[pages=-]{reports/schematic.pdf}
\section{Layout}
Gate \texttt{place} (phase P6): \textbf{pass} after 31 attempt(s); last 2026-09-27T08:55:02, failing 0 of 0.

Gate \texttt{drc\_routed} (phase P7): \textbf{pass} after 14 attempt(s); last 2026-09-27T08:55:04, failing 0 of 0.
\begin{center}
\includegraphics[width=0.72\textwidth]{reports/PCB-0017-A_stereo-class-d-amp_top.png}\\
{\small\texttt{reports/PCB-0017-A\_stereo-class-d-amp\_top.png}}
\end{center}
\begin{center}
\includegraphics[width=0.72\textwidth]{reports/PCB-0017-A_stereo-class-d-amp_bottom.png}\\
{\small\texttt{reports/PCB-0017-A\_stereo-class-d-amp\_bottom.png}}
\end{center}
\section{Board Layers}
Each of the board's 2 copper layers as KiCad's editor shows it, seen from the top: the copper with its side's silkscreen, fabrication and courtyard outlines, and each pad, track and zone labelled with its net. The pages are vector drawings, so zooming in keeps the names sharp.
\clearpage
\subsection*{F.Cu}
\begin{center}
\includegraphics[width=\textwidth,height=0.82\textheight,keepaspectratio]{reports/layers/PCB-0017-A_stereo-class-d-amp_F_Cu.pdf}\\
{\small\texttt{reports/layers/PCB-0017-A\_stereo-class-d-amp\_F\_Cu.pdf}}
\end{center}
\clearpage
\subsection*{B.Cu}
\begin{center}
\includegraphics[width=\textwidth,height=0.82\textheight,keepaspectratio]{reports/layers/PCB-0017-A_stereo-class-d-amp_B_Cu.pdf}\\
{\small\texttt{reports/layers/PCB-0017-A\_stereo-class-d-amp\_B\_Cu.pdf}}
\end{center}
\clearpage
\subsection*{3D views}
\paragraph*{top}
\begin{center}
\includegraphics[width=0.8\textwidth]{reports/layers/PCB-0017-A_stereo-class-d-amp_top.png}\\
{\small\texttt{reports/layers/PCB-0017-A\_stereo-class-d-amp\_top.png}}
\end{center}
\paragraph*{bottom}
\begin{center}
\includegraphics[width=0.8\textwidth]{reports/layers/PCB-0017-A_stereo-class-d-amp_bottom.png}\\
{\small\texttt{reports/layers/PCB-0017-A\_stereo-class-d-amp\_bottom.png}}
\end{center}
\paragraph*{iso}
\begin{center}
\includegraphics[width=0.8\textwidth]{reports/layers/PCB-0017-A_stereo-class-d-amp_iso.png}\\
{\small\texttt{reports/layers/PCB-0017-A\_stereo-class-d-amp\_iso.png}}
\end{center}
\section{Routing Renders}
The finished board's renders are under Layout; this section shows how it got there.
\subsection*{Snapshots the run kept}
The run saved 17 board snapshots before (pre-) or after (post-) a step that changed the board; they are shown in the order the run made them.
 None is rendered yet: report\_gen.py --kind full --render-history renders them.
\section{Design Decisions}
Every decision the run recorded in state.json (37), oldest first, with the reason given at the time.
{\small
\begin{longtable}{lp{1.6cm}p{6.2cm}p{6.2cm}}
\hstoprule
\hshead{Ph} & \hshead{When} & \hshead{Decision} & \hshead{Why} \\
\hline
\endhead
P0 & 09-27 00:36 & Q1 pads: J1=12V in (V+,GND), J2=stereo line in (L,R,GND), J3=left spk +/-, J4=right spk +/-; TP1 per the reference design & delegate answer (owner said decide yourself); matches the reference design's J1-J4 \\
P0 & 09-27 00:36 & Q2 speaker: 4 ohm min, 8 ohm typ & delegate; 10uH filter suits both \\
P0 & 09-27 00:36 & Q3 power in: bare wire pads & delegate; like-for-like with the reference design \\
P0 & 09-27 00:36 & Q4 no mute/fault LED; fixed pull-ups, always on & delegate; the reference design has none \\
P0 & 09-27 00:36 & Q5 supply 12V nom, 10.8-13.2V, caps \textgreater{}=25V & delegate default \\
P0 & 09-27 00:36 & Q6 PROVISIONAL (safety, \textgreater{}3A borderline): 3A continuous on 12V input, 3A per speaker output & delegate; TPA3116D2 \textasciitilde{}2x15W into 4ohm at 12V -\textgreater{} \textasciitilde{}2.8A supply; unattended run, reconfirm at H1 \\
P0 & 09-27 00:36 & Q7 no reverse-polarity protection & delegate; matches the reference design part count; noted as a deviation risk in COMPARISON \\
P0 & 09-27 00:36 & Q8 PCB copper heat-sinking only (PowerPAD thermal vias to bottom pour), grow board if needed and record & owner: thermal first, then size \\
P0 & 09-27 00:36 & Q9 no mounting holes & delegate; bare module like the reference design \\
P0 & 09-27 00:36 & Q10 qty 5, no unit-cost target, JLC stocked parts & delegate default \\
P0 & 09-27 00:36 & Q11 JLC PCBA top side only unless size forces otherwise & delegate default \\
P0 & 09-27 00:36 & Q12 line-level in, keep the reference design's 5.6k/100k gain/config & delegate; like-for-like \\
P1 & 09-27 00:53 & DEVIATION: amp IC TPA3116D2DAD -\textgreater{} TPA3118D2DAP (same SLOS708 family, same 32-pin HTSSOP pinout, PowerPAD DOWN) & TPA3116D2 exists only as DAD (pad UP, TI SLOS708); its 14 C/W needs a top heatsink, which the reference design's own write-up admits is not modeled. Brief: thermal pad/heat-sinking first, SMD only. At 12 V output is rail-limited so no audio power lost; pad-down lets 2-layer copper + vias sink it (\textasciitilde{}22-30 C/W). \\
P1 & 09-27 00:53 & Q6 revised: +12V input and its GND return sized for 4 A; 3 A per speaker output; inductor Isat \textgreater{}= 4 A & power research: 3.6 A at 10\% THD into 4 ohm, 4.0 A at 13.2 V \\
P1 & 09-27 00:53 & Sustained full-power 4 ohm sine at 40 C is NOT a requirement (music + 8 ohm full power are); keep the size target & power research: that case needs \textasciitilde{}50x50 mm 2 oz or PLIMIT; thermal shutdown is the chip's own protection; FAULTZ-\textgreater{}SDZ gives auto-recovery \\
P1 & 09-27 00:53 & Prefer 2 oz outer copper if the stackup table offers it for 2L & halves trace widths and improves spreading; 0.65 mm pitch is fine at 2 oz floors \\
P1 & 09-27 00:53 & PLIMIT tied to GVDD (no limit); FAULTZ tied to SDZ with 100k pull-up for self-clearing thermal fault; gain 20 dB via 5.6k master & TI 12 V design; refdesign research resolves the reference design's 5.6k/100k as gain leg + pull-ups \\
P1 & 09-27 00:53 & No reverse-polarity/TVS part (Q7 stands) & like-for-like with the reference design; risk (reversed supply, hot-plug ring) written in COMPARISON.md \\
P1 & 09-27 01:13 & 10 uF PVCC decoupling in 0805 (Basic CL21A106KAYNNNE), not 1206 & Basic part, smaller footprint helps size; 25 V rating kept \\
P1 & 09-27 01:13 & No switch-node EMI snubber (10R+330pF); keep only the 3.3R+10nF damping branch per channel output & like-for-like with the reference design and the EVM damping network; EMI not a stated requirement \\
P1 & 09-27 01:13 & Parts leads: TPA3118D2DAPR C46497; SLO0630H100MTT 10uH C207842 (Isat 5.5A); 680nF out filter; bulk MA25V100M6x6 polymer + RVT 100uF 35V electrolytic & P1 component-scout shortlist \\
P2 & 09-27 01:39 & P2 research tiers: class-d-amp + lc-output-filter researchers on fable/high (tier); dc-input + audio-line-in researchers and second readers on sonnet/high (explicit downgrade) & B1/B2 are low-risk textbook blocks already covered by P1 power.md/refdesign; owner asked for quality without burning money \\
P2 & 09-27 01:42 & Start P3 part-sourcer in parallel with P2 research, before H1 & part choice doesn't turn on the checklist research; a value the research changes is amended in parts.json before P4 \\
P3 & 09-27 02:00 & C11/C12 10 uF PVCC: accept X5R CL21A106KAYNNNE (C15850) though sheets.md said X7R & caps sit below 85 C (X5R limit) at \textless{}=1.5 W design dissipation; DC-bias loss at 12 V is covered by the 100 uF bulk per pair; Basic part \\
P3 & 09-27 02:00 & R4-R7 3.3R 0603 100 mW (RC0603FR-073R3L) accepted & branch dissipation \textasciitilde{}1 mW: 10 nF is \textasciitilde{}800 ohm at 20 kHz and \textasciitilde{}40 ohm at 400 kHz against \textless{}0.3 V residual ripple \\
P2 & 09-27 02:06 & R4-R7+C23-C26 (3.3R+10nF) stay on the speaker side of the LC filter as TI's EMI C-RC snubber, not at the switch node & research: TI labels it an EMI snubber (corner \textasciitilde{}4.8 MHz), not audio damping; at the switch node 10 nF x 12 V\textasciicircum{}2 x 400 kHz = 0.58 W would overload a 0603. Filter Q 0.52-1.04 over 4-8 ohm is acceptable per SLAA701B \\
P3 & 09-27 02:56 & designing under coverage gap: block:B1 dc-input classes power-loop, esd, sourcing, constraints-emission have no record at the approved floor; no second research round & B1 is two wire pads + two 100uF bulk caps with no protection by decision Q7; its loop/placement rules come from the verified B3 PVCC three-tier decoupling and dc-input hot-plug records; a 2nd round would cost \textasciitilde{}\$5-8 for little design delta \\
P3 & 09-27 02:56 & no TI Fig 37 PVCC RC damper (C58 100nF + R10 3.3R) added & C14 electrolytic ESR is the bulk damping per the verified dc-input-bulk-esr-damping-budget record; hot-plug ring margin 1.15x stays a P4-review watch item with RC snubber as fallback \\
P5 & 09-27 03:27 & U1 footprint fix: the 9 EasyEDA thermal vias (0.3 drill, 0.61 pad, 1.0 pitch 3x3) had empty pad numbers (no net -\textgreater{} 9 mask-bridge + 9 clearance errors); renumbered to pad 33 (GND) and set layers {*}.Cu only (no paste/mask on the vias: bottom tented, top inside the EP mask window) & thermal vias must be GND and join the bottom pour; record class-d-amp-powerpad-thermal-via-array asks component-side capping, which is impossible inside the EP window without plugging - via-in-pad plugging is the P9 alternative; 3x3 at 1.0 mm meets the \textgreater{}=6 count / \textless{}=1.2 mm pitch rule \\
P5 & 09-27 03:27 & wire-pad footprints: moved pin-1 silk corner 0.35 mm outside pad 1 (was over copper, 4 silk\_over\_copper warnings) & footprint-internal silk is a librarian fix; board\_init self-check now 0 violations \\
P4 & 09-27 03:36 & review W1/W2: swap C13 to a \textgreater{}=35 V low-ESR polymer/hybrid D6.3 part and C14 to a 35 V part with a stated ESR/ripple; ADD a unidirectional TVS (SMBJ15A-class, clamp \textless{} 30 V abs max) across J1 +12V/GND & reviewer: 26 V hot-plug ring exceeds C13's 25 V and C14's ESR cannot damp the bank; TVS bounds the ring below U1's 30 V abs max. Correctness over the reference design's part count (deviation for COMPARISON) \\
P4 & 09-27 03:36 & review W3: C1-C4 2.2 uF -\textgreater{} 1 uF X7R 50 V 0805, K tolerance (TI Fig 37 value; corner 2.7 Hz at 60k) & 2.2 uF x 60k (132 ms) with no source connected reaches the 420 ms DC-detect window at tolerance corners and fault-cycles on power-up; deviation from the reference design's 2.2 uF recorded \\
P4 & 09-27 03:36 & review W4: reverse polarity stays unprotected (Q7) but J1 gets a +12V/GND silk legend; W5: C11/C12 X5R 10 uF accepted & W4 accepted risk, legend is free; W5: \textasciitilde{}4-5 uF effective at 12 V is ample beside 200 uF bulk + 100 nF/1 nF HF caps \\
P4 & 09-27 03:46 & Schematic review W1-W3 fixed: C13 35V polymer C46550467, C14 FZ 35V C2992596 (stated ESR/ripple), C1-C4 1uF X7R C28323, D1 SMBJ15A C113988 across J1 & W1 26V ring \textgreater{} 25V C13 and C14 cannot damp it -\textgreater{} TVS clamps; W2 RVT had no ESR spec; W3 132ms input tau vs 420ms DC-detect -\textgreater{} 60ms \\
P6 & 09-27 04:35 & Accept 53.6 x 36.8 mm outline (2 layers) as a deviation from the 44 x 34 target & Thermal/decoupling layout first: 25x25 mm clear thermal copper around U1, PVCC stacks at the pins, 7.3 mm inductors need 9.6 mm column pitch and 3x6 mm speaker wire pads J3/J4 take their own column; closing \textasciitilde{}9 mm would need smaller inductors or wire pads, trading current rating/solderability for size \\
P6 & 09-27 04:35 & Strip body-outline silk from C0402/C0603/C0805/R0603 (lib + board) and widen IND courtyard to cover pads & 98 outline-vs-outline silk\_overlap warnings in the dense PVCC/bootstrap stacks would fail drc\_routed; passives are unpolarised, refdes kept; IND courtyard was smaller than its pads \\
P7 & 09-27 05:25 & P7-\textgreater{}P6: spread L3/L4 up and L1/L2 down \textasciitilde{}2.5 mm (grow outline height if needed) to fit the C15-C18 bootstrap caps in two columns and route OUT nets on F.Cu & OUT nets unroutable with current placement; caps-on-B.Cu alternative would cut the U1 heat-spreader pour edge with 8 vias; thermal first, area already a stated deviation \\
\hline
\end{longtable}
}
\section{Design History}
\subsection*{Phase timeline}
{\small
\begin{longtable}{llllr}
\hstoprule
\hshead{Phase} & \hshead{Name} & \hshead{First decision} & \hshead{Last decision} & \hshead{Decisions} \\
\hline
\endhead
P0 & Intake & 2026-09-27T00:36 & 2026-09-27T00:36 & 12 \\
P1 & Research & 2026-09-27T00:53 & 2026-09-27T01:13 & 9 \\
P2 & Architecture & 2026-09-27T01:39 & 2026-09-27T02:06 & 3 \\
P3 & Parts and library & 2026-09-27T02:00 & 2026-09-27T02:56 & 4 \\
P4 & Schematic & 2026-09-27T03:36 & 2026-09-27T03:46 & 4 \\
P5 & Board setup & 2026-09-27T03:27 & 2026-09-27T03:27 & 2 \\
P6 & Placement & 2026-09-27T04:35 & 2026-09-27T04:35 & 2 \\
P7 & Routing & 2026-09-27T05:25 & 2026-09-27T05:25 & 1 \\
\hline
\end{longtable}
}
\subsection*{Human checkpoints}
{\small
\begin{longtable}{lllp{9cm}}
\hstoprule
\hshead{H} & \hshead{Status} & \hshead{When} & \hshead{Note} \\
\hline
\endhead
1 & approved & 2026-09-27T03:07 & delegate (owner: decide for yourself). Coverage: C46497 covered; B2/B3/B4 provisional, all 4 research tasks closed verified; B1 gap accepted by decision. TPA3118D2 DAP pad-down instead of 3116 DAD pad-up (needs heatsink); 2L 2oz; 4 A input; riskiest: 2L thermal at 4 ohm sustained, hot-plug ring margin 1.15x. \\
2 & approved & 2026-09-27T03:46 & delegate (owner: decide for yourself). Schematic 44 parts/29 nets, ERC 0/0, netlist\_audit pass; review 0 err/5 warn, W1-W3 fixed, W4 silk legend at layout, W5 X5R accepted. Deviation from the reference design: TPA3118D2 pad-down + SMBJ15A TVS added + 2x100uF bulk. \\
3 & approved & 2026-09-27T04:35 & delegate (owner: decide for yourself). Place gate PASS, DRC 0 after silk fix, 53.6x36.8 mm 2L accepted as deviation (decision). \\
\hline
\end{longtable}
}
\subsection*{Where the run went back}
\paragraph*{P3 back to P2, 2026-09-27T02:06}
\begin{itemize}
\item R4-R7+C23-C26 (3.3R+10nF) stay on the speaker side of the LC filter as TI's EMI C-RC snubber, not at the switch node \emph{Why:} research: TI labels it an EMI snubber (corner \textasciitilde{}4.8 MHz), not audio damping; at the switch node 10 nF x 12 V\textasciicircum{}2 x 400 kHz = 0.58 W would overload a 0603. Filter Q 0.52-1.04 over 4-8 ohm is acceptable per SLAA701B
\end{itemize}
\paragraph*{P5 back to P4, 2026-09-27T03:36}
\begin{itemize}
\item review W1/W2: swap C13 to a \textgreater{}=35 V low-ESR polymer/hybrid D6.3 part and C14 to a 35 V part with a stated ESR/ripple; ADD a unidirectional TVS (SMBJ15A-class, clamp \textless{} 30 V abs max) across J1 +12V/GND \emph{Why:} reviewer: 26 V hot-plug ring exceeds C13's 25 V and C14's ESR cannot damp the bank; TVS bounds the ring below U1's 30 V abs max. Correctness over the reference design's part count (deviation for COMPARISON)
\item review W3: C1-C4 2.2 uF -\textgreater{} 1 uF X7R 50 V 0805, K tolerance (TI Fig 37 value; corner 2.7 Hz at 60k) \emph{Why:} 2.2 uF x 60k (132 ms) with no source connected reaches the 420 ms DC-detect window at tolerance corners and fault-cycles on power-up; deviation from the reference design's 2.2 uF recorded
\item review W4: reverse polarity stays unprotected (Q7) but J1 gets a +12V/GND silk legend; W5: C11/C12 X5R 10 uF accepted \emph{Why:} W4 accepted risk, legend is free; W5: \textasciitilde{}4-5 uF effective at 12 V is ample beside 200 uF bulk + 100 nF/1 nF HF caps
\item Schematic review W1-W3 fixed: C13 35V polymer C46550467, C14 FZ 35V C2992596 (stated ESR/ripple), C1-C4 1uF X7R C28323, D1 SMBJ15A C113988 across J1 \emph{Why:} W1 26V ring \textgreater{} 25V C13 and C14 cannot damp it -\textgreater{} TVS clamps; W2 RVT had no ESR spec; W3 132ms input tau vs 420ms DC-detect -\textgreater{} 60ms
\end{itemize}
\paragraph*{P7 back to P6, 2026-09-27T05:25}
\begin{itemize}
\item P7-\textgreater{}P6: spread L3/L4 up and L1/L2 down \textasciitilde{}2.5 mm (grow outline height if needed) to fit the C15-C18 bootstrap caps in two columns and route OUT nets on F.Cu \emph{Why:} OUT nets unroutable with current placement; caps-on-B.Cu alternative would cut the U1 heat-spreader pour edge with 8 vias; thermal first, area already a stated deviation
\end{itemize}
\subsection*{Component choices and changes}
{\small
\begin{longtable}{lp{1.6cm}p{6.2cm}p{6.2cm}}
\hstoprule
\hshead{Ph} & \hshead{When} & \hshead{Decision} & \hshead{Why} \\
\hline
\endhead
P1 & 09-27 00:53 & DEVIATION: amp IC TPA3116D2DAD -\textgreater{} TPA3118D2DAP (same SLOS708 family, same 32-pin HTSSOP pinout, PowerPAD DOWN) & TPA3116D2 exists only as DAD (pad UP, TI SLOS708); its 14 C/W needs a top heatsink, which the reference design's own write-up admits is not modeled. Brief: thermal pad/heat-sinking first, SMD only. At 12 V output is rail-limited so no audio power lost; pad-down lets 2-layer copper + vias sink it (\textasciitilde{}22-30 C/W). \\
P1 & 09-27 01:13 & Parts leads: TPA3118D2DAPR C46497; SLO0630H100MTT 10uH C207842 (Isat 5.5A); 680nF out filter; bulk MA25V100M6x6 polymer + RVT 100uF 35V electrolytic & P1 component-scout shortlist \\
P3 & 09-27 02:00 & C11/C12 10 uF PVCC: accept X5R CL21A106KAYNNNE (C15850) though sheets.md said X7R & caps sit below 85 C (X5R limit) at \textless{}=1.5 W design dissipation; DC-bias loss at 12 V is covered by the 100 uF bulk per pair; Basic part \\
P5 & 09-27 03:27 & U1 footprint fix: the 9 EasyEDA thermal vias (0.3 drill, 0.61 pad, 1.0 pitch 3x3) had empty pad numbers (no net -\textgreater{} 9 mask-bridge + 9 clearance errors); renumbered to pad 33 (GND) and set layers {*}.Cu only (no paste/mask on the vias: bottom tented, top inside the EP mask window) & thermal vias must be GND and join the bottom pour; record class-d-amp-powerpad-thermal-via-array asks component-side capping, which is impossible inside the EP window without plugging - via-in-pad plugging is the P9 alternative; 3x3 at 1.0 mm meets the \textgreater{}=6 count / \textless{}=1.2 mm pitch rule \\
P4 & 09-27 03:36 & review W1/W2: swap C13 to a \textgreater{}=35 V low-ESR polymer/hybrid D6.3 part and C14 to a 35 V part with a stated ESR/ripple; ADD a unidirectional TVS (SMBJ15A-class, clamp \textless{} 30 V abs max) across J1 +12V/GND & reviewer: 26 V hot-plug ring exceeds C13's 25 V and C14's ESR cannot damp the bank; TVS bounds the ring below U1's 30 V abs max. Correctness over the reference design's part count (deviation for COMPARISON) \\
P4 & 09-27 03:36 & review W3: C1-C4 2.2 uF -\textgreater{} 1 uF X7R 50 V 0805, K tolerance (TI Fig 37 value; corner 2.7 Hz at 60k) & 2.2 uF x 60k (132 ms) with no source connected reaches the 420 ms DC-detect window at tolerance corners and fault-cycles on power-up; deviation from the reference design's 2.2 uF recorded \\
P4 & 09-27 03:46 & Schematic review W1-W3 fixed: C13 35V polymer C46550467, C14 FZ 35V C2992596 (stated ESR/ripple), C1-C4 1uF X7R C28323, D1 SMBJ15A C113988 across J1 & W1 26V ring \textgreater{} 25V C13 and C14 cannot damp it -\textgreater{} TVS clamps; W2 RVT had no ESR spec; W3 132ms input tau vs 420ms DC-detect -\textgreater{} 60ms \\
P6 & 09-27 04:35 & Accept 53.6 x 36.8 mm outline (2 layers) as a deviation from the 44 x 34 target & Thermal/decoupling layout first: 25x25 mm clear thermal copper around U1, PVCC stacks at the pins, 7.3 mm inductors need 9.6 mm column pitch and 3x6 mm speaker wire pads J3/J4 take their own column; closing \textasciitilde{}9 mm would need smaller inductors or wire pads, trading current rating/solderability for size \\
P6 & 09-27 04:35 & Strip body-outline silk from C0402/C0603/C0805/R0603 (lib + board) and widen IND courtyard to cover pads & 98 outline-vs-outline silk\_overlap warnings in the dense PVCC/bootstrap stacks would fail drc\_routed; passives are unpolarised, refdes kept; IND courtyard was smaller than its pads \\
P7 & 09-27 05:25 & P7-\textgreater{}P6: spread L3/L4 up and L1/L2 down \textasciitilde{}2.5 mm (grow outline height if needed) to fit the C15-C18 bootstrap caps in two columns and route OUT nets on F.Cu & OUT nets unroutable with current placement; caps-on-B.Cu alternative would cut the U1 heat-spreader pour edge with 8 vias; thermal first, area already a stated deviation \\
\hline
\end{longtable}
}
\subsection*{Git history}
1 commits touched the workspace, 2026-09-27T14:53 to 2026-09-27T14:53. By hour, with the files each area had changed:
{\small
\begin{longtable}{lrp{10cm}}
\hstoprule
\hshead{Hour} & \hshead{Commits} & \hshead{Areas (files)} \\
\hline
\endhead
2026-09-27 14:00 & 1 & routing 154, research 69, lib 37, reports 36, fab 19, anneal 17, kicad 16, route 15, log 11, (root) 5, architecture 5, parts 2, brief 1, learnings 1 \\
\hline
\end{longtable}
}
\subsection*{Comparison}
\subsection*{PCB-0017-A stereo-class-d-amp vs a published reference design}
Same brief as a published 44 x 34 mm reference design: TPA311x stereo class-D amp on one 12 V supply, stereo in and out, SMD only, 2 layers, as small as possible. The reference design's numbers are from JLCPCB's write-up (https://jlcpcb.com/blog/gpt-6-astra-pcb-design-in-kicad ; research/refdesign-amp.md, "the reference design's board as reported").

\begin{quote}\small\ttfamily\raggedright
\textbar{} \textbar{} Reference (reported) \textbar{} hwde, PCB-0017 rev A \textbar{}\\
\textbar{}---\textbar{}---\textbar{}---\textbar{}\\
\textbar{} Board size \textbar{} 44 x 34 mm (1496 mm2) \textbar{} 53.6 x 36.8 mm (1972 mm2, +32 \%) \textbar{}\\
\textbar{} Layers \textbar{} 2 \textbar{} 2 (1.6 mm, 2 oz) \textbar{}\\
\textbar{} Footprints \textbar{} 45 \textbar{} 44 \textbar{}\\
\textbar{} Fitted (BOM) parts \textbar{} 40 \textbar{} 39 (J1-J4 and TP1 are bare wire pads, as in the reference design's) \textbar{}\\
\textbar{} Amplifier \textbar{} TPA3116D2 \textbar{} TPA3118D2 (same die family, thermal pad down) \textbar{}\\
\textbar{} ERC \textbar{} 0 \textbar{} 0 errors, 0 warnings \textbar{}\\
\textbar{} DRC \textbar{} 0 violations, 0 unconnected \textbar{} 0 violations, 0 unconnected (all severities) \textbar{}\\
\textbar{} JLCPCB DFM \textbar{} not reported \textbar{} pass, 0 findings \textbar{}\\
\textbar{} Deeper checks (current, return path, thermal) \textbar{} not reported \textbar{} 8 findings, all proposed as waivers for the owner (below) \textbar{}\\
\textbar{} Fabricated / measured \textbar{} no \textbar{} no (fab package only, nothing ordered) \textbar{}
\end{quote}
\subsubsection*{Where we deviate, and why}
\begin{itemize}
\item \textbf{Bigger board.} We put heat and decoupling first. The pad-down TPA3118D2 dumps its heat through
the exposed pad into the copper around it, so I kept a clear 25 x 25 mm GND area around U1 on both
layers and the PVCC decouplers tight to their pins. The 7.3 mm, 10 uH inductors need a 9.6 mm column
pitch, and the 3 x 6 mm speaker wire pads take their own column. Closing the last \textasciitilde{}9 mm would mean
smaller inductors (lower saturation current) or smaller wire pads, so I left it.
\item \textbf{TPA3118D2, not TPA3116D2.} The TPA3116D2 only ships pad-up (DAD), which needs a heatsink bolted
on top. The TPA3118D2 is the pad-down (DAP) sibling and sinks into the board, which is what a
small 2-layer SMD-only board needs.
\item \textbf{Added D1, an SMBJ15A TVS across the 12 V input,} and rated the bulk caps at 35 V (polymer C13,
low-ESR C14) for hot-plug ringing. The reference design's list has no input protection.
\item \textbf{1 uF input caps} instead of 2.2 uF: against the chip's input impedance the high-pass corner
stays at a few hertz, well below the audio band, and the part is a cheaper basic X7R.
\item The reference design's damping branch (3.3 ohm + 10 nF per channel) and 680 nF filter caps are kept.
\end{itemize}
\subsubsection*{What neither board has proven}
Neither board has been built. The reference design's write-up reported DRC/ERC-clean and nothing more. Our verify pass goes further and flags eight things: the thermal-copper check, whose 2-layer model can't show 1.5 W fitting under 85 C on any pour, six pad-width links at U1's 0.65 mm pins that fall below the IPC-2152 current width, and one GND pour neck on F.Cu. I think all eight are artefacts of the checks rather than defects, and each is written up as a proposed waiver in reports/verify-waivers.proposed.json for the owner to accept or refuse. The thermal case rests on TI's measured 22 C/W on its 2-layer EVM; a fresh-context review of our layout put the first layout at 35-40 C/W, mostly because U1's pad had only nine thermal vias. U1 now has a 4 x 4 array of sixteen, which I estimate brings it to about 27-32 C/W, so the chip sits near 85 C at the 1.5 W design point and 40 C ambient, inside its 125 C limit. The review also asked for bootstrap caps C16 and C17 to sit closer to U1, and they now sit about 4.7 mm from their pins instead of 9. To make room, C15 moved out to about 4.5 mm from its pin, and the 1 nF and 100 nF supply caps C7 and C9 moved just south of U1. The real test of both boards is a sustained 4-ohm run on the bench.

\subsection*{Placement notes}
\subsection*{PCB-0017-A stereo-class-d-amp - P6 placement notes (2026-09-27)}
\textbf{Chosen: a hand-built canonical placement (anneal/manual3-5.ops.json), not the seed or an anneal candidate.} The seed (anneal/seed.ops.json, HPWL 2187 -\textgreater{} 856 mm) was legal but spread over the 74x97 provisional outline; the annealer cannot shrink an outline, so the compact layout was written as absolute ops around the datasheet rules.

Layout (all top side, U1 rot 270: pins 1-16 face left/J2, pins 17-32 face right):

\begin{itemize}
\item J2 line-in on the left edge, R/GND/L top to bottom; C3/C4/C1/C2 input caps and R2/C5/R1/R3 between J2 and U1's pin 1-16 side.
\item PVCC 31/32 (top right): C8 1nF, C10 100nF, C12 10uF in a stack right at the pins, C13 100uF above them, D1 and J1 (+12V left pad, GND right pad) at the top edge.
\item PVCC 18/19 (bottom right): C7, C9, C11, C6 (AVCC) stack, C14 100uF below.
\item Bootstrap C18/C17/C16/C15 vertical in two staggered columns beside their BS/OUT pin pairs.
\item L3/L4 (right ch, top row) and L1/L2 (left ch, bottom row) as a 2x2 grid right of U1, OUT pads facing U1; a 4.6 mm channel between rows carries OUTNR/OUTNL to the second column and holds R7/C26/C22/C20.
\item Right-channel 680nF + Zobel (C21, R6, C25) above the grid; left-channel (C19, R4, C23, R5, C24) below it. J4 (upper) and J3 (lower) on the right edge, pads long-axis horizontal.
\item TP1 bottom-left. Nothing sits under U1; the 25x25 thermal square holds only U1, its decouplers and the input parts.
\end{itemize}
Gate: place PASS (0 violations) after 5 edit iterations (manual1..5). silk\_place --apply moved 44 refdes, 0 residuals; gate still passes. Outline: board\_edit --outline fit --margin 0.8 -\textgreater{} bbox {[}9.1, 8.3, 62.65, 45.1{]} = 53.6 x 36.8 mm (vs 44 x 34 reference).

Why larger than 44 x 34: (1) the inductor footprint's pads (2.5 x 3.1 at +/-3.08) reach +/-4.34 mm, outside its 6.7 x 6.6 courtyard, so the 2x2 grid needs a 9.6 mm column pitch; (2) the gate's courtyard model is pads + \textasciitilde{}0.5 mm, so every stack needs 0.5-0.6 mm gaps; (3) J3/J4 (7.2 x 6.8 wire pads) sit outside the second inductor column instead of overlapping it vertically; (4) the bulk caps are 8.8 mm wide pad-to-pad. Recovering \textasciitilde{}8 mm width means stacking J3/J4 between the inductor rows and the board corners, or swapping the 7.3 mm inductors for a 6 mm part.

Not done in the 28 min cap: route\_auto --probe (routability), full kicad-cli DRC after silk, "PCB-0017 REV A" and "+12V/GND" silk texts. Room for them: the REV text fits along the bottom edge between TP1 and C14 (x 19-31, y 43.5); "+12V/GND" fits above J1's pads at y 9.2 (x 44-51). Open items: IND-SMD\_L7.3-W6.6\_MHCI06030 courtyard smaller than its pads (footprint fix in lib); bootstrap-cap OUT traces must leave the cap pads sideways (see anneal/manual3.ops.json layout); SPK\_L\_P route to J3's lower pad goes under J3's upper pad, so route it at y \textasciitilde{}43 from x 46. Backups: kicad/PCB-0017-A\_stereo-class-d-amp.kicad\_pcb.p5bak (board\_init state), anneal/pass-before-silk.kicad\_pcb.

\subsection*{Routing notes}
\subsection*{PCB-0017-A stereo-class-d-amp - P7 routing notes (router subagent, 2026-09-27 08:37-09:02Z)}
Hard time cap 25 min. Snapshots before every mutating step in routing/pre-\textless{}step\textgreater{}.kicad\_pcb.

\subsubsection*{Chain run}
1. planes\_gen: PASS. Two GND pours added (B.Cu 1878 mm2, F.Cu 1558 mm2, both one island), 9 thermal vias at C14/D1/J1 GND pads. U1's EP vias were already present. Constraint fix needed first: \texttt{planes{[}{]}.\{\_why\}} is rejected by the validator; the two notes were moved to a top-level \texttt{\_planes\_why} key (pre-edit copy: routing/constraints.pre-router.json). 2. route\_critical --pad-window: PASS, every +12V pad ceiling \textgreater{}= 0.882 mm (widest 5.4-8.0 mm). route\_critical (power, KRT) was NOT run to completion: the first KRT invocation exceeded a 60 s cap and the default is 600 s per net, which did not fit the window. Skipped by decision; the netclasses (Pwr\_0p624mm / Pwr\_0p882mm) carry the IPC widths so Freerouting honours them. 3. route\_auto: two aborted starts before a live run - a) DISPLAY=:99 had no Xvfb (only :163/:187 exist) -\textgreater{} export\_dsn failed; used :163. b) Freerouting's bundled JRE could not load libXtst.so.6 -\textgreater{} both rungs died at JVM start. Fix: LD\_LIBRARY\_PATH=\$HOME/.local/kicad10/usr/lib/x86\_64-linux-gnu (KiCad bundle ships it). hwde-env.sh should export this itself (tooling finding, not a board issue). Live run: --max-rungs 2 --timeout-s 150 --no-krt-finish. Result: see below.

\subsubsection*{Knowledge records applied}
\begin{itemize}
\item class-d-amp-2layer-pour-heat-spreader: B.Cu GND pour is the heat spreader; kept as one island.
\item audio-line-in-2layer-switch-node-keepaway: no shield layer on 2L, so line-in and OUT nets are
separated by placement (J2 left, filters right); routing must not cross them on B.Cu under U1.
\item lc-output-filter-placement-and-loop / class-d-amp-pvcc-three-tier-decoupling: decoupler returns
go straight to the pour (F.Cu GND pour connect solid per constraints.planes).
\end{itemize}
\subsubsection*{Result (08:52Z, stopped at the time cap)}
\begin{itemize}
\item Board = planes\_gen state + nothing else. Freerouting never produced a SES: both rungs
(150 s each) were killed by the per-rung timeout while still iterating. Its log shows it was
close: pass 1 = 6 unrouted (81 s), passes 2-4 = 5 unrouted (\textasciitilde{}20 s each, score flat at 976.18),
i.e. 5 nets it cannot finish with this placement - likely a placement\_adjust candidate, but
route\_auto never reached the point of emitting that request.
\item Copper: 0 tracks, 0 signal vias; GND pours F/B (one island each) + EP thermal vias + 9 pad vias.
\item Gate drc\_routed: FAIL, 56 unconnected (was 90 before pours), 0 other DRC errors/warnings.
Report: reports/gate-drc\_routed.json. Not recorded in state.json (a failing gate).
\item check\_thermal: not run (no routing to screen yet).
\item Nothing hand-fixed. No placement change requested yet.
\end{itemize}
\subsubsection*{Exact next step (P7 continues)}
Same shell env plus: export DISPLAY=:163   (or start Xvfb :99 first; it was not running) export LD\_LIBRARY\_PATH=\$HOME/.local/kicad10/usr/lib/x86\_64-linux-gnu 1. Optional but wanted for the +12V trunk / OUT widths: route\_critical.py --pcb kicad/... (default 600 s KRT timeout; budget \textasciitilde{}10-15 min). Skipped here for time. 2. route\_auto.py --pcb kicad/PCB-0017-A\_stereo-class-d-amp.kicad\_pcb --work-dir route --out-report routing/route\_auto.json with --timeout-s \textgreater{}= 600 (Freerouting needs to finish a run to write the SES; 150 s was not enough). Read facts.completion; if it carries placement\_adjust\_request for the \textasciitilde{}5 stubborn nets, return it to the orchestrator (P7-\textgreater{}P6 backward edge). 3. stitch\_vias -\textgreater{} plane\_repair (snapshot first) -\textgreater{} route\_cleanup --dry-run then live -\textgreater{} gate. 4. check\_thermal.py --pcb ... --constraints kicad/constraints.json, then state.py record-gate. Tooling findings to feed back to hwde: hwde-env.sh should export the KiCad lib dir for the JRE (libXtst.so.6) and either start or verify the Xvfb display it expects; planes\_gen's validator should accept \texttt{\_}-prefixed annotation keys inside planes{[}{]} like the rest of constraints.json does.

\subsection*{P7 routing, second router (2026-09-27 09:03-09:20Z, hard stop 09:36Z)}
Start: route\_auto2 result = 170 tracks, 16 DRC (8 track\_width, 3 edge clearance, 5 unconnected). Snapshots: routing/pre-route\_edit\_fix.kicad\_pcb, pre-refill, pre-rulearea; end state copied to routing/post-router2.kicad\_pcb.

\subsubsection*{Steps}
1. route\_edit (routing/route\_edit\_fix.json, 22 ops, report route\_edit\_fix.report.json):

\begin{itemize}
\item +12V left-edge run moved from x=9.806 to x=9.9 (3 tracks re-added, 0.882 mm) -\textgreater{} edge clearance
0.36 mm, fixes the 3 copper\_edge\_clearance errors.
\begin{itemize}
\item +12V (31.5,34.6)-(29.9,35.109) three 0.80 mm segments re-added at 0.882 mm.
\item Removed Freerouting's /OUTPR route: it went through a via at (27.886,24.018) UNDER U1's body
to B.Cu, ran diagonally under the heat-spreader pour to (34.96,16.875) and detoured to L3 from
above. Wrong per the class-d heat-spreader and outputs-on-TOP records, and 0.343 mm wide.
\end{itemize}
\end{itemize}
2. Zones refilled (kicad-cli pcb drc --refill-zones --save-board); the stale fill was hugging the old thin tracks and gave 4 clearance errors against the GND pour after step 1. 3. Rule exception, pad-entry neck-down at U1 (0.65 mm pitch, 0.343 mm pads; a 0.882/0.624 mm track physically cannot enter these pads - max is 0.345 mm for 0.1524 mm clearance):

\begin{itemize}
\item Rule area "U1\_neck" on F.Cu, x 28.6-30.95, y 21.6-32.4 (U1's right pad column + 0.35 mm),
no keepouts (pure named area). Added with the bundled pcbnew python (route\_edit has no op for it).
\begin{itemize}
\item kicad\_dru: rule "aiee\_neck\_fine\_pitch\_U1", track\_width min 0.1524 mm, condition
A.Type=='track' \&\& A.intersectsArea('U1\_neck'), appended LAST so it wins inside the area.
A.intersectsCourtyard('U1') and A.insideCourtyard('U1') did NOT match in KiCad 10 here
(courtyard shapes exist) - tooling note. rules\_gen.py would overwrite this rule; pre-edit
copy at routing/pre-neck-rule.kicad\_dru. The two remaining 0.343 mm tracks are pin-to-pin
bridges between adjacent +12V pads (17-18-19, 31-32) inside the pad column.
\end{itemize}
\end{itemize}
4. plane\_repair --flag-only: PASS, no split planes (B.Cu pour still one piece). route\_cleanup --dry-run: 0 ops, so no live run. stitch\_vias NOT run: routing is unfinished and stitch vias would land in the region the OUT re-route needs. 5. gate drc\_routed: FAIL 6/6 - 0 DRC errors, 6 unconnected (4 OUT nets). reports/gate-drc\_routed.json. Not recorded (failing). 6. check\_thermal -\textgreater{} reports/check\_thermal.json: 1 error + 1 warning on U1 (1.5 W, "0 thermal vias", \textasciitilde{}111 C rise). It counts board vias only; U1's 9 thermal vias are footprint pads with drills (pad 33, 0.61 mm) so it sees none. Tooling finding; the 645 mm2 copper figure is also suspect given the 1878 mm2 B.Cu pour. Re-run after the OUT routing is done before acting on it.

\subsubsection*{Why the four OUT nets are still open (needs a placement change, P7 -\textgreater{} P6)}
U1 right column pins, top to bottom: 30 BSPR 23.42, 29 OUTPR 24.07, 28 GND, 27 OUTNR 25.37, 26 BSNR 26.02, 25 GND, 24 BSPL 27.33, 23 OUTPL 27.98, 22 GND, 21 OUTNL 29.28, 20 BSNL 29.93. Bootstrap caps today: near column x=32.2 (C18: BSPR 23.4 / OUTPR 25.4; C16: BSPL 27.6 / OUTPL 29.6), far column x=34.2-34.4 (C17: OUTNR 25.5 / BSNR 27.5; C15: OUTNL 29.5 / BSNL 31.5).

\begin{itemize}
\item C18's OUTPR pad (y 24.7-26.1) sits exactly in the path of OUTNR (pin y 25.37) to C17; C16's
OUTPL pad (y 28.9-30.3) sits in the path of OUTNL (29.28) to C15. Both blocked on F.Cu.
\item The gap between the columns is 0.65 mm; a 0.624 mm OUT track needs 0.93 mm, so OUTPR cannot
climb from C18 to L3 between the caps either - that is why Freerouting dived under U1.
\item Space between the PVCC pin groups is 7.6 mm for four 3.4 mm-tall caps in two columns plus
eight tracks; a re-place of C15-C18 is required, not a nudge. Worked (right channel only):
C18 -\textgreater{} (32.8,23.4) rot -90 (clears the +12V run at x=31.5/y=21.7), C17 -\textgreater{} far column \textasciitilde{}(38.4,26.55)
rot 90, OUTNR row y=25.55, BSNR row y=26.15, BSPL 27.15. The left channel does not close in the
same scheme: OUTPL must drop to L1 (y 33.6) past the OUTNL/BSNL rows whichever column C16/C15
take. Candidate: move L1 up/left or put C15/C16 above/below the pin band and re-run
place gate + route\_auto with C15-C18 fixed; or route OUTPL/OUTNL as F.Cu polygons.
\end{itemize}
\subsubsection*{Files changed by this router}
kicad/PCB-0017-A\_stereo-class-d-amp.kicad\_pcb (tracks, rule area, zone refill), kicad/PCB-0017-A\_stereo-class-d-amp.kicad\_dru (neck rule), routing/route\_edit\_fix.json + .report.json, routing/pre-{*}.kicad\_pcb, routing/pre-neck-rule.kicad\_dru, routing/post-router2.kicad\_pcb, routing/plane\_repair\_flag.json, routing/route\_cleanup\_dry.json, reports/gate-drc\_routed.json, reports/check\_thermal.json, this file. Nothing committed, nothing ordered.

\subsubsection*{09:22-09:25Z: P7-\textgreater{}P6 attempt, stopped before touching the board (coordinator cap 09:44Z)}
Region survey (x 30.5-48.5, y 19.5-37) shows the re-place is bigger than a cap nudge:

\begin{itemize}
\item C8 (32.2,21.7) bbox y 21.25-22.15 and C7 (32.2,33.1) bbox y 32.65-33.55 bracket the near column:
usable y for caps at x\textasciitilde{}32.8 is 22.3-32.5.
\item L3 body bbox reaches y 25.86 and L1 body bbox starts at y 30.15, both from x 35.9 (pads only at
y 20.8-24.0 / 32.0-35.2). So right of x=35.9 there are only 4.3 mm of free height (25.9-30.1)
for the far-column caps; two 3.35 mm-tall caps do not stack there, and the rows for the inner
pairs (OUTNR 26.14, BSNR 26.70, BSPL 27.06, OUTPL 27.62 - derived from C8's edge downward) leave
no cap-sized gap either.
\item Existing BS tracks to remove when re-routing: /BSPR 39c81251, 911ab399, 9227e97e; /BSNR 2d981dc7,
49ccdffe, 98385004, ac5aa048, ae3ec111, e5955bab; /BSPL 1d8b337c, c05d49bd, c6908c3d;
/BSNL 023ccccb, 08417efb, 561f827f, e927c9e3. /OUTNL and /OUTNR partials from C15/C17 to L2/L4 exist.
\end{itemize}
What closes it (pick one in P6, then re-run place gate, route\_auto or hand route): a) Shift L3/L4 (+SPK\_R tracks, R7/C26/ZR\_N cluster) up \textasciitilde{}2.5 mm and L1/L2 (+SPK\_L) down \textasciitilde{}2.5 mm, or grow the outline; then the two-column cap arrangement above fits (C18 pads y 23.0/25.0 at x 32.8; C15 pads 28.77/30.77 at x 32.8; C17 and C16 in a far column x\textasciitilde{}38.4). b) Put C15-C18 on B.Cu under x 31.5-35 (east of U1's EP, pour stays one piece), free F.Cu for four straight OUT tracks U1 -\textgreater{} L1..L4, tie the caps with vias. Cheaper on area, costs 8 vias in the spreader edge. Board left exactly at routing/post-router2.kicad\_pcb: 0 DRC errors, 6 unconnected, outline unchanged (53.65 x 36.9 mm, 9.05-62.7 x 8.25-45.15). Place gate not re-run (no placement change).

\subsubsection*{09:24-09:32Z: option (a) started, stopped at the fallback (coordinator cap 09:45Z)}
Done on the board, then saved as routing/wip-optionA.kicad\_pcb and the clean board restored (kicad/...kicad\_pcb == routing/post-router2.kicad\_pcb, drc\_routed re-run: 0 DRC errors, 6 unconnected): 1. route\_edit rip (routing/optionA\_rip.json, 44 removed): all /BSPR /BSNR /BSPL /BSNL tracks, the /OUTNL and /OUTNR partials (C15-\textgreater{}L2, C17-\textgreater{}L4), all /SPK\_L\_P and /SPK\_L\_N tracks. /ZL\_P, /ZL\_N, /ZR\_{*} and /SPK\_R\_{*} were NOT ripped (right channel not moved). ZL\_{*} stubs at R4/R5/C23/C24 will dangle in the wip and must be ripped when resuming. 2. place\_edit move (routing/optionA\_move.json, --allow-routed because other nets stay routed): L1 (40.5,36.4), L2 (50.1,36.4), C19 (40.5,41.4), R4 (43.9,41.4), C23 (46.3,41.4), R5 (50.5,40.7), C24 (50.5,42.3) - the left channel down 2.8 mm (2.5 was 0.03 mm short for two caps in the far column). Outline NOT changed (53.65 x 36.9 mm); C24's bottom lands at y 42.6 vs edge 45.15. 3. gate place on the wip: FAIL 5 - C11/L1 courtyard 0.10 mm2 and C14/L1 0.60 mm2 (new: L1 cannot shift right, L2 and J3 are hard against it, so C11 and C14 have to give way instead, \textasciitilde{}0.5 mm each), C12/L3 0.21 mm2 (pre-existing), and U1 + C8 "intersect rule area U1\_neck" (2). reports/gate-place-optionA.json. The right channel was not moved: above L3/L4 sit C21/R6/C25 (y 17.4) and then J1/D1 at the top edge, so "up 2.5 mm" means growing the outline upward by \textasciitilde{}2.5 mm and moving J1, D1, C13, C12, C21, R6, C25, L3, L4, J4 plus the +12V trunk. That is the next pass's first decision; the left half of option (a) is ready in the wip.

Place gate on the CLEAN board (reports/gate-place.json): FAIL 4 - C11/L1 0.08 mm2 and C12/L3 0.21 mm2 courtyard overlaps (pre-existing, never recorded), plus U1 and C8 intersecting the U1\_neck area. The place check treats any rule area as a keepout; U1\_neck has no keepouts and only exists as the DRC condition for the fine-pitch neck rule, so this is a tooling conflict between check.place and the neck exception. gate.py --waivers is for human-approved findings, so I left it alone. Ways out for the next pass: teach check.place to skip rule areas with no keepout flags, or find a KiCad-10 expression that matches the two 0.343 mm pad-to-pad +12V bridges without an area. Neither gate recorded (both failing). Nothing committed or ordered.

\subsection*{09:36-10:00Z: place gate closed, neck rule area gone; OUT nets still open (hard stop 10:12Z)}
Board left at routing/post-placefix.kicad\_pcb (== kicad/...kicad\_pcb) + post-placefix.kicad\_dru. Pre-state snapshots: routing/pre-placefix.kicad\_pcb/.kicad\_dru (== post-router2), pre-placefix2.

\subsubsection*{Result}
\begin{itemize}
\item gate place: PASS, recorded in state.json (attempt 18). reports/gate-place.json.
\item gate drc\_routed: FAIL 6/6 - 0 DRC errors, 0 warnings, 6 unconnected (the 4 OUT nets, as before).
Not recorded. GND pours: F.Cu 1 island, B.Cu 1 island (EP via array untouched).
\item Outline unchanged: 53.55 x 36.8 mm edge rect (9.1-62.65 x 8.299-45.1). Silk "PCB-0017 REV A",
"+12V", "GND" untouched.
\end{itemize}
\subsubsection*{What moved (place\_edit --allow-routed, files routing/placefix{*}\_move.json / \_route.json)}
\begin{itemize}
\item C11 -\textgreater{} (34.6, 35.6) rot 90 (vertical, +12V pad below at 36.6, GND above at 34.6). It had 0.08 mm2
courtyard overlap with L1; horizontally there is no room between C6's GND pad (x 32.05) and L1's
courtyard (x 35.905) for an 0805 pad field (3.35 mm) plus check.place's margin, so it went vertical.
+12V feed re-routed: (32.629,38.089)-\textgreater{}(34.6,38.089)-\textgreater{}(34.6,36.6), 0.882 mm.
\item C12 -\textgreater{} (34.65, 19.2) rot -90 (+12V pad on top at 18.2). It overlapped L3 by 0.21 mm2 and at any x
between the trunk vertical (x 33.6) and L3 it either shorted the trunk or overlapped L3. The +12V
trunk vertical moved left 0.3: (35.05,15.118)-\textgreater{}(33.3,16.868)-\textgreater{}(33.3,17.834)-\textgreater{}(32.026,19.108), plus a
stub (33.3,18.2)-\textgreater{}(34.65,18.2) into the pad. All 0.882 mm.
\item Reference texts moved (they collided with L1/L3 silk and pads after the moves): C11 ref (26.6,33.9),
C9 ref (34.7,33.3), C12 ref (36.9,17.4). move\_text takes field "reference" (lower case).
\item check.place's effective courtyard = courtyard UNION pad box, plus a margin of roughly 0.3 mm - the
KiCad courtyards on these 0805s are body-only (2.09 x 1.33) while the pad field is 3.35 x 1.41.
Plan positions against pad boxes, not courtyards.
\end{itemize}
\subsubsection*{U1\_neck resolved}
The rule area is deleted (text edit of the .kicad\_pcb; a pcbnew SWIG Save() was tried first and reverted - not the cause, but avoided). The dru rule "aiee\_neck\_fine\_pitch\_U1" now uses A.Type == 'track' \&\& A.NetName == '+12V' \&\& A.Start\_X \textgreater{} 29mm \&\& A.Start\_X \textless{} 30.5mm \&\& A.End\_X \textgreater{} 29mm \&\& A.End\_X \textless{} 30.5mm i.e. an x-window on U1's right pad column (pads span x 28.855-30.585), which matches exactly the two 0.343 mm pad-to-pad bridges (pins 17-18-19 and 31-32). Verified with kicad-cli drc: 0 track\_width. Tested and rejected (kicad-cli, /tmp copies): \texttt{A.Length \textless{} 1.0mm} does NOT parse in KiCad 10 and silently disables the WHOLE dru (symptom: netclass clearances only, pours refill with 0.5 mm clearance, F.Cu GND splits into 9 islands, 17 unconnected). \texttt{A.intersectsCourtyard('U1')} parses but never matches because U1's courtyard is body-only (x 22.955-29.045) and excludes the pads. If rules\_gen.py is re-run, re-append this rule from routing/post-placefix.kicad\_dru. Tooling findings: (1) check.place should ignore rule areas with no keepout flags; (2) a dru parse failure is silent in kicad-cli - gate.py should fail loudly when the dru does not load (detect by running with and without the dru and comparing the rule-derived violation counts, or grep the log).

\subsubsection*{Not done: option (a) re-place of C15-C18 and the OUT/BS/SPK\_L routing}
Nothing of option (a) was applied (the 25 min went to the place gate + neck rule). Worked geometry for the next pass, on top of THIS board (C11/C12 already fixed), all F.Cu, clearance 0.2 to power:

\begin{itemize}
\item Pad facts: bootstrap caps are 0805, pads at +-1.0 from centre, pad box 1.35 x 1.41, so a vertical
cap occupies y c-1.705..c+1.705. C8 pad bottom 22.15, C7 pad top 32.65 bound the near column.
\item Near column x=32.8: C18 rot -90 centre y 24.1 (BSPR 23.1, OUTPR 25.1); C15 rot 90 centre 30.2
(OUTNL 29.2, BSNL 31.2). Four rows pass between them: OUTNR 26.35 (0.624), BSNR 26.96 (0.2),
BSPL 27.36 (0.2), OUTPL 27.98 (0.624) - the clearances are 0.2-0.235, so keep these numbers.
\item Far column x=38.4 (needs L1/L2 down 2.8 as in wip-optionA: L1 (40.5,36.4), L2 (50.1,36.4), C19
(40.5,41.4), R4 (43.9,41.4), C23 (46.3,41.4), R5 (50.5,40.7), C24 (50.5,42.3); L3/L4 stay):
C17 rot 90 centre 27.35 (OUTNR 26.35, BSNR 28.35); C16 rot -90 centre 31.05 (BSPL 30.05,
OUTPL 32.05). Drops in x order: OUTPL down at x=36.2 (27.98 -\textgreater{} 34.2, then diagonal into L1 pad 1
at (37.42,35.4); stub (36.2,32.05)-\textgreater{}(38.4,32.05)); BSPL down at x=36.9 (27.36 -\textgreater{} 30.05, then right);
BSNR down at x=37.4 (26.96 -\textgreater{} 28.35, then right); OUTNR straight into C17 at 26.35, then on to
(44.6,26.35)-\textgreater{}(47.0,23.95)-\textgreater{}(47.0,22.6) = L4 pad 1. OUTPR: (32.8,25.1)-\textgreater{}(35.9,25.1)-\textgreater{}(37.42,23.6)
= L3 pad 1.
\item OUTNL is the hard one: it must cross OUTPL's drop. Take it down WEST of OUTPL: from C15's pad right
to x=35.0, down to y 39.0, east under L1's pads (pad bottom 37.96 after the move) to x=45.55, up to
36.4 and into L2 pad 1 (x 46.258+). That x=35.0 vertical needs C11 clear of x 34.69-35.31: C11 as
placed now (pad box 33.925-35.275) is 0.03 too far right -\textgreater{} shift C11 to x=34.3 first (pad box left
33.625, still 0.3 clear of C9's GND pad at 33.3). C14 must also move left \textasciitilde{}0.7 (its courtyard hits
the moved L1: 0.60 mm2 in the wip) and its +12V track (30.23,39.757-\textgreater{}41.0) re-added.
\item Inside U1's pad column use 0.3 mm stubs to x=30.8 before widening to the class width (the GND
pads 22/25/28 are 0.2 away from the OUT pads; a 0.624 mm end cap at x\textless{}30.8 violates). The neck
rule only covers +12V, so add a second x-window rule for the OUT nets (or widen the net list) before
routing them, else 4 track\_width errors.
\item Rip first (routing/optionA\_rip.json, 44 uuids, still valid on this board): all BS{*}, the OUTNL/OUTNR
partials, SPK\_L\_\emph{; also rip ZL\_} (4+1 tracks) since R4/R5/C23/C24 move.
\item Then: place gate -\textgreater{} hand route (route\_edit) or route\_auto with --timeout-s 480 --no-krt-finish -\textgreater{}
kicad-cli drc --refill-zones --save-board -\textgreater{} stitch\_vias -\textgreater{} route\_cleanup -\textgreater{} gate drc\_routed -\textgreater{}
record. Budget: 45-60 min; it did not fit in this window.
\end{itemize}
\subsection*{10:00-10:05Z: option (a) applied once, did not converge, fallback taken (hard stop 10:18Z)}
Board restored to routing/post-placefix.kicad\_pcb (+ .kicad\_dru); place gate re-run PASS and re-recorded; drc\_routed unchanged (0/0/6). The attempt is saved whole as routing/wip-optionA3.kicad\_pcb, wip-optionA3.kicad\_dru, wip-optionA3.drc.json, with its op lists routing/optionA3\_move.json (13 place ops) and optionA3\_route.json (52 rips + 71 adds, all applied and verified by route\_edit).

What the attempt found (all on the wip, one pass, no iteration):

\begin{itemize}
\item DRC copper: 4 clearance errors, 1 short, 1 mask bridge, 1 unconnected; the F.Cu GND pour split
into 2 islands (B.Cu still 1). Causes: (1) BSNR's start (30.9,26.02)-\textgreater{}(30.9,26.26) sits 0.15 from
OUTNR's 45-degree rise - start BSNR's rise further right (x 31.2) instead of a vertical jog;
(2) OUTNL's x=35.4 vertical and its y=40.0 run hit C14's GND pad (C14 at 31.2 puts pad 2 at 33.87
with a \textasciitilde{}2 mm pad: right edge \textasciitilde{}34.9, top \textasciitilde{}40.0) - either run OUTNL's horizontal at y 39.4 with L1
at y 37.4 (pad bottom 38.96, too tight) or put C14 at x\textless{}=30.6; (3) SPK\_L\_P's y=40.95 run from the
via to C19 crosses C19's GND pad (41.5,41.75) - route it at y 40.4 or enter C19 pad 1 from the
left; the mask bridge is the same track. The pour split is the OUTNL loop (x 35.4 down, y 40 east,
x 45.55 up) enclosing the L1 area: a GND island inside a ring of OUT copper - add a stitch via
inside the ring (far from U1) or leave a gap for the pour.
\item Place gate: 5 courtyard overlaps - check.place's margin is \textasciitilde{}0.35-0.4 mm around pad boxes, not
0.2-0.3: C18/C8 (0.245 gap), C17/L3 (crt 25.695 vs pad top 25.645), C16/C17 (0.2 gap),
C11/C9 (same-net GND pads 0.175 apart still count), C11/C7. So the two-column stack needs \textasciitilde{}0.4 mm
more per gap: C18 centre 24.3, rows from 26.55, C15 centre \textasciitilde{}30.6 (C7 pad top 32.65 -\textgreater{} bottom
limit 32.3, ok), far column C17 centre 27.6, C16 centre 31.5 -\textgreater{} L1 down to \textasciitilde{}37.8, then C19/R4/
C23/R5/C24 and the SPK\_L\_P bottom run at y \textasciitilde{}44.0 (edge 45.1) - or grow the outline 1 mm down.
\item Silk: refs of C16, C18, C24, C9 need moving after the re-place; "PCB-0017 REV A" was moved to
(18,12) in the wip because C14 at x=31.2 reaches it at (24.5,43.6).
\item The dru neck rule in the wip drops the NetName clause and widens the window to x 29-31 mm so
the 0.3/0.2 mm OUT/BS stubs (29.72 -\textgreater{} 30.9) pass; keep that when resuming.
\end{itemize}
Estimated to close: one more 45-60 min pass from wip-optionA3 fixing the seven items above.

\subsection*{10:06-10:28Z: option (a) converged - all nets routed, place PASS, drc\_routed PASS (hard stop 10:45Z)}
Live board = kicad/PCB-0017-A\_stereo-class-d-amp.kicad\_pcb + .kicad\_dru == routing/wip-optionA4.kicad\_pcb/.kicad\_dru (dru unchanged from wip-optionA3: neck rule with no NetName clause, window x 29-31 mm). Resumed from wip-optionA3; op lists routing/optionA4\_rip.json (57 removes: every OUT\emph{/BS}/SPK\_L\_P track+via plus the two +12V C11 feed tracks), optionA4\_route.json (65 adds), optionA4\_route2.json (2 GND stitch vias), optionA4\_move.json (+ move2..move5: C14 and reference-text retries). Reports alongside ({*}.report.json).

\subsubsection*{Gate results (both recorded in state.json)}
\begin{itemize}
\item gate place: PASS 0/0 (attempt 23). reports/gate-place.json.
\item gate drc\_routed: PASS 0/0 (attempt 7): kicad-cli DRC --severity-all 0 violations (silk included),
0 unconnected. reports/gate-drc\_routed.json. Zones refilled (--refill-zones --save-board).
\item GND pours: B.Cu 1 polygon. F.Cu 4 polygons, all connected (DRC unconnected 0): the U1-side region
west of the OUT/BS rows, the C7/C9/C11 pocket between U1's column and OUTNL's x=35.5 drop, and the
region south of OUTNL's y=40 run are each tied to B.Cu by vias (U1's EP array; new stitch vias at
(34.3,34.9) and (37.0,42.3); plus (41.0,34.6) inside the OUTNL/OUTPL ring). Not literally "one
island" - that is impossible on F.Cu once the four OUT nets run from U1 to the far column.
\end{itemize}
\subsubsection*{What changed vs wip-optionA3 (all F.Cu, OUT 0.624, BS 0.2, stubs 0.3/0.2 at x 29.72-30.9)}
\begin{itemize}
\item The key move: C18 is now HORIZONTAL (rot 0) at (33.8, 23.36) - BSPR pad (32.8), OUTPR pad (34.8).
This takes OUTPR out of the between-column stack, and the vertical budget between C18 and C15
(which the place gate's 0.25 mm-per-pad-box margin made 0.06 mm too short for 4 rows) is no longer
needed there. Rows between the columns: OUTPR 24.59 (passes under C18 pad 1, into pad 2, then
y 23.36 straight into L3 pad 1), OUTNR 25.42, BSNR 26.12, BSPL 26.53 (rises from its pin), OUTPL
27.98 -\textgreater{} 28.4 after clearing C15's pad box at x 33.6.
\item Far column: C17 (38.4, 27.7) rot 90 (OUTNR 26.7 / BSNR 28.7); C16 (38.4, 31.65) rot -90 (BSPL
30.65 / OUTPL 32.65). Near column: C15 (32.8, 30.4) rot 90 (OUTNL 29.4 / BSNL 31.4).
Crossings east of x 34: OUTNR diag (35.0,25.42)-(36.28,26.7); BSNR diag (34.5,26.12)-(35.72,27.34)
then row 27.34 -\textgreater{} drop x 37.42 -\textgreater{} 28.7; BSPL diag (33.5,26.53)-(34.72,27.75) then row 27.75 -\textgreater{}
drop x 37.0 -\textgreater{} 30.65. OUTPL drop x 36.33 -\textgreater{} 32.65 -\textgreater{} C16 pad -\textgreater{} (38.4,35.9) straight into L1 pad 1.
\item OUTNL: row 29.25 -\textgreater{} x 35.5 down to y 40.0 -\textgreater{} east to 45.55 -\textgreater{} up to L2 pad 1 (as A3 but 0.1 right).
\item L1, L2, C19, R4, C23, R5, C24, L3, L4 stay where wip-optionA3 put them; outline NOT grown
(53.55 x 36.8 unchanged) - not needed once C18 went horizontal.
\item C11 -\textgreater{} (34.15, 37.3) rot 90 (GND 36.3 / +12V 38.3): 0.5 mm from C9's pad box (place gate counts
same-net pads), +12V feed is the single track (32.629,38.089)-(34.15,38.089) ending in the pad.
\item C14 -\textgreater{} (29.8, 41.0): its pad 2 (3.5 x 1.2 mm!) now ends at x 34.22, clear of OUTNL; its round
courtyard reaches x 33.35 so C14 at 30.2 still overlapped C11 (0.43 mm2). The +12V track into
pad 1 (ends 28.53,41.0) is unchanged and still lands in the pad.
\item SPK\_L\_P: L1 pad 2 -\textgreater{} via (43.58,39.0) -\textgreater{} B.Cu -\textgreater{} via (43.58,40.85) -\textgreater{} F.Cu (43.15,41.28) -\textgreater{} R4
pad 1 -\textgreater{} down to y 43.95 -\textgreater{} J3 as before; C19 pad 1 is fed by a branch from (43.15,42.97) running
BELOW C19's GND pad at y 42.97 and entering pad 1 from below (fixes the short + mask bridge).
\item Silk: refs moved - C18 (33.8,25.05), C16 (40.3,31.65), C11 (34.6,39.75), C24 (50.2,44.35).
"PCB-0017 REV A" at (18,12) untouched. Full-severity DRC: 0 silk warnings.
\item route\_edit gotcha: a "remove" and an "add\_track" of identical geometry in ONE op list rolls back
("not in saved board") - the add is skipped as "exists" against the pre-removal board. Rip and
add in two runs.
\end{itemize}
\subsubsection*{P8 verify\_all preview (/tmp/a4verify.json, not recorded): 9 errors, 21 warnings}
\begin{itemize}
\item Known / owner waiver, not mine: check\_thermal thermal\_area U1 (2L model); 2x undersized\_track
+12V 0.343 mm at x 29.72 (U1 PVCC pin-to-pin bridges).
\item Same class as those two, needs the same waiver: 4x undersized\_track /OUTPR /OUTNR /OUTPL /OUTNL
0.300 mm at (30.31, pin y) - the 1.18 mm pad-entry stubs inside the neck window; a 0.65 mm pitch
pad cannot take a 0.624 mm track.
\item check\_current pour\_neckdown GND F.Cu at (56.24, 26.70) "\textasciitilde{}0.00 mm between via attachments" - NEW
location (the earlier one near (45.67,30.76) is gone); next to SPK\_L\_N's x=56.0 vertical. Not
looked at; likely a stitch via whose pour attachment is pinched between the SPK\_L\_N track and
L4/J4 copper. Move or drop that via.
\item check\_current insufficient\_transition\_vias /SPK\_L\_P at (43.58,39.9): 2 vias, 3 A wants 6. Real
finding from the layer hop. Either accept for a 1-2 A speaker load (owner call) or add vias: there
is no room for 6 between L1 pad 2 (bottom 38.96) and OUTNL (top 39.688); via-in-pad on L1 pad 2
plus a row below OUTNL (y \textasciitilde{}40.85, x 42.0-44.0, above R4/C19 pads at 41.3) fits \textasciitilde{}3+3.
\item Warnings: 5 GND transition-via counts (advisory), 7 check\_decoupling (C5/C8/C10/C12/C13 GND stubs
and loops - pre-existing, not touched), 7 check\_silk.
\end{itemize}
stitch\_vias.py and route\_cleanup.py were NOT run this pass (no time); run both before P8 for real. Nothing committed, nothing ordered.

\subsection*{10:28-10:36Z: SPK\_L\_P layer hop fixed (12 vias); pour\_neckdown is structural, not fixed}
Live board == routing/wip-optionA5.kicad\_pcb/.kicad\_dru (dru unchanged). Op lists routing/optionA5\_rip.json, optionA5\_route.json, optionA5\_rip2.json, optionA5\_route2.json (+ reports). Gates re-recorded: place PASS 0/0, drc\_routed PASS 0/0 (full-severity DRC 0 violations, 0 unconnected, zones refilled).

\begin{itemize}
\item SPK\_L\_P: check\_current counts vias per END of a transition (6 each at 3 A), so: 6 via-in-pad on
L1 pad 2 (x 42.7/43.55/44.4 at y 37.75 and 38.6), B.Cu (43.55,38.6)-\textgreater{}(43.55,43.5)-\textgreater{}(44.0,43.95)-\textgreater{}
(49.55,43.95), 6 vias on the F.Cu y=43.95 run at x 45.3..49.55 (0.85 pitch). OUTNL's bottom run
raised to y 39.9. The 2-via hop and the 40.75 bus are gone. insufficient\_transition\_vias: cleared.
\item pour\_neckdown GND F.Cu "(56.24,26.70)": NOT a real neck. pour\_neck() erodes the whole F.Cu GND
fill by 0.441 mm and wants every GND via in one eroded component; the F.Cu GND is 4 polygons
(OUT rows cut the U1 pocket off, OUTNL's loop cuts two more), each with vias (U1 EP array, the
stitch vias), so the width is 0 and the reported point is just parts{[}0{]}'s representative point
(C20's GND pad). Two GND vias I added at (57.2,27.0)/(54.8,27.0) did nothing; left in (harmless).
Fix is either a check that works per polygon, or a waiver alongside thermal\_area. Do not chase it
by moving SPK\_L\_N.
\item verify\_all preview (/tmp/a5verify.json): 8 errors = 2x +12V 0.343 (waiver), 4x OUT 0.300 stubs
(waiver, same class), 1 pour\_neckdown (above), 1 check\_thermal thermal\_area (waiver). 23 warnings
(transition-via advisories incl. the new via rows, check\_decoupling C5/C8/C10/C12/C13, check\_silk).
\end{itemize}
\subsection*{10:36-10:38Z: SPK\_L\_P via field moved off L1 pad 2 (no via-in-pad)}
Live board == routing/wip-optionA6.kicad\_pcb/.kicad\_dru (dru unchanged). Ops routing/optionA6\_rip.json (11), optionA6\_route.json (13). The 6 vias are now under L1's body between its pads, x 39.6/40.45/41.3 at y 36.9 and 37.75 (0.85 pitch; 0.6/0.3), fed from pad 2 by two 0.624 F.Cu rows (39.6..43.0) and joined on B.Cu, then B.Cu (41.3,37.75)-\textgreater{}(41.3,43.5)-\textgreater{}(41.75,43.95)-\textgreater{}(44.0,43.95) into the existing 6-via row on the y=43.95 run. No via sits inside any SMD pad. Gates re-recorded: place PASS 0/0, drc\_routed PASS 0/0 (full-severity DRC 0/0). verify\_all preview: 8 errors, the same 8 (2x +12V + 4x OUT undersized pin-entry stubs, pour\_neckdown structural, thermal\_area); 23 warnings; no SPK\_L\_P transition finding.

\subsection*{11:45-12:00Z: EP 4x4 array + C16/C17 bootstrap placement (hard stop 12:20Z)}
Fallback = routing/pre-epfix.kicad\_pcb/.kicad\_dru (the live board before this pass). Live board now = routing/epfix\_post-ep.kicad\_pcb (dru unchanged). Review items 2 and 4 of reports/review-verify.md.

\subsubsection*{Fix 1 (done): U1 exposed-pad thermal vias 3x3 @1.0 -\textgreater{} 4x4 @0.9 mm}
\begin{itemize}
\item 16 pad-"33" thru-holes at x,y in \{-1.35,-0.45,0.45,1.35\} (footprint-local), 0.61 pad / 0.30 drill,
outer via edge 1.655 from centre inside the 3.81 x 3.71 EP. Library footprint and the board's U1
instance (at (26,27) rot -90, all other pads/nets untouched) rewritten by routing/epfix\_ep\_pads.py;
lib/EDITS.md has the entry. No track or via sat under the EP on either layer, so nothing was ripped.
\item Refill + DRC: routing/epfix\_ep\_refill.drc.json = 0 violations. Both GND zones are connect\_pads yes,
so all 16 vias are solid-connected to the B.Cu pour and sit inside the F.Cu EP copper.
\item Gates recorded (state.json + reports/gate-\emph{.json; copies routing/epfix\_gate-}.json): place PASS 0/0
(attempt 27), drc\_routed PASS 0/0 full severity, 0 unconnected (attempt 11), dfm PASS 0/0 (attempt 4).
\item fp\_verify (routing/epfix\_fp\_verify.json): pad\_count 49 != 32 error is the checker counting the EP +
thermal vias as copper pads (baseline gave 42 != 32 plus a pad\_pitch error that is now gone);
pad\_size warning unchanged. Not a footprint fault.
\item verify\_all preview (routing/epfix\_verify.json, checks in /tmp/epfix\_checks): 8 errors, SAME set,
same positions: +12V undersized 0.343 at (29.72,30.905) and (29.72,22.445); OUT 0.300 stubs at
(30.31, 27.98/29.28/24.07/25.37) for OUTPL/OUTNL/OUTPR/OUTNR; pour\_neckdown GND F.Cu at
(56.2389,26.6984); check\_thermal thermal\_area U1 at (26.0,27.0). 23 warnings, unchanged.
Note: check\_thermal's thermal\_vias warning still says "found 0 via(s), want \textgreater{}= 15" - it counts board
vias only, not the footprint's pad-33 thru-holes, before and after; tooling gap, not a board fact.
\end{itemize}
\subsubsection*{Fix 2 (NOT done, board unchanged for it): C16/C17 within \textasciitilde{}4 mm of pins 24/26}
C16 stays (38.4,31.65) rot -90 = 9.4 mm from pin 24 (29.72,27.33); C17 stays (38.4,27.7) rot 90 = 8.8 mm from pin 26 (29.72,26.02). Worked out and rejected in the 20 min available:

\begin{itemize}
\item The only free pocket within 4 mm of the pins is the slot between the OUTPR row (bottom edge 24.90)
and C15's pad box top (28.695), 3.79 mm tall; a vertical 0805 needs 3.41 + 2x0.5 place margin, so
only HORIZONTAL caps fit, and two of them need 0.2+1.35+0.5+1.35+0.5 = 3.9 mm (0.11 short; C7, C9,
C15 down 0.2 would give it - their +12V track ends still land inside the 0.9 mm pads).
\item But a horizontal cap forces the OUT net to pass the BS pad on F.Cu (OUT pad on the left: the
onward OUT track to L1/L4 must pass the BS pad; on the right: the pin track must). Above is the
OUTPR row, below is the other cap: no 0.624 + 2x0.2 corridor. The 3 A OUT nets cannot hop to B.Cu
(check\_current wants 6 vias per transition end).
\item What would work (next pass, 45-60 min like the option-(a) passes): C15 horizontal too, three
horizontal caps at x centre 33.0 (pads 32.0/34.0) stacked at y \textasciitilde{}25.8 (C17, rot 180: OUTNR pad left),
\textasciitilde{}28.2 (C16, rot 180: OUTPL pad left) and \textasciitilde{}30.6 (C15), each OUT net passing the cap above it in a
0.624 corridor (cap pitch 2.375), BS rows at 0.2 threading between (BSPL row \textasciitilde{}27.6 needs C17 bottom
\textless{}= 27.28 and C16 top \textgreater{}= 27.9). Requires re-routing OUTNR/BSNR/OUTPL/BSPL/OUTNL/BSNL from the
pins to the existing tracks east of x 38.4, moving the OUTNL x=35.5 drop, and re-doing C15's nets;
L1-L4 stay. Distances would be \textasciitilde{}2.3-3.0 mm to the BS pins.
\end{itemize}
Nothing committed, nothing ordered.

\subsection*{11:58-12:12Z: C16/C17 bootstrap placement - reverted, geometry re-worked (hard stop 12:27Z)}
Live board = routing/pre-bsfix.kicad\_pcb/.kicad\_dru (== the 4x4-EP board from the pass above), restored at 12:07Z; gate place PASS 0/0 and drc\_routed PASS 0/0 re-run on it (reports/gate-{*}.json, copies routing/bsfix\_gate-{*}.json). C15/C16/C17 are where they were. Nothing committed, nothing ordered. Geometry dump of the region: routing/bsfix\_geometry.txt (routing/bsfix\_dump.py, bundled pcbnew python: \textasciitilde{}/.local/kicad10/bin/python3 - the venv has no pcbnew, the bundled python has no shapely).

\subsubsection*{The plan above (three horizontal caps at x 33) cannot be routed on one layer}
Pin order top to bottom on U1's east side: OUTPR 24.07, BSNR-pair: OUTNR 25.37 / BSNR 26.02; BSPL 27.33 / OUTPL 27.98; OUTNL 29.28 / BSNL 29.93. For a horizontal 0805 both nets of a pair arrive from the left; whichever pad is on the right, the net going to it fences the cap off from the pin column, so the OTHER net's onward track (the OUT net must go on east to L1-L4) has to pass on the far side of the cap. With the OUT pin ABOVE the BS pin (OUTNR/BSNR, OUTNL/BSNL) that far side is the top: OUT pad left =\textgreater{} OUT onward must go over the cap; BS pad left =\textgreater{} OUT must cross BS's direct pin-to-pad track. Above C17 sits the OUTPR row (bottom edge 24.90, and it cannot rise: C18's pads end at 24.035), so there is no 0.624+2x0.2 corridor: C17 horizontal is dead in that slot, and so is C15 horizontal below C16 (the corridor above it would be under the OUTPL row). Only BS-above-OUT pairs (BSPR/OUTPR = C18 today, BSPL/OUTPL = C16) work horizontal with BS left.

\subsubsection*{What does work: vertical caps, pads in pin order (nothing crosses)}
\begin{itemize}
\item C17 VERTICAL at (33.0, 26.85) rot 90: pad 2 OUTNR top (33.0,25.85), pad 1 BSNR bottom (33.0,27.85);
pad box x 32.325-33.675, y 25.145-28.555 (0.245 under the OUTPR row). 3.4 mm from pin 26.
Applied and DRC'd in routing/wip-bsfix.kicad\_pcb (ops bsfix\_rip\_c17 / bsfix\_move\_c17 /
bsfix\_route\_c17 .json + reports; DRC routing/bsfix\_c17.drc.json, 14 violations): all but one are
the NOT-yet-ripped BSPL 26.53 row and OUTPL 27.98/28.4 row and C15 pad 2 hitting the new cap -
expected, they go with C16/C15. The one real one: OUTNR (30.9,25.37)-\textgreater{}(31.2,25.6) is 0.104 from
the BSNR stub; start the diagonal later, e.g. (30.9,25.37)-\textgreater{}(31.6,25.6)-\textgreater{}(32.6,25.6), and enter the
top pad from its left; onward east (33.4,25.85)-\textgreater{}(35.0,25.85)-\textgreater{}(36.28,26.7) joins the untouched
OUTNR track to L4 (10 tracks ripped: BSNR east of x 30.95, OUTNR 31.2..36.28).
\item C16 VERTICAL at (33.0, 30.76) rot 90 with pad 1 BSPL on TOP (29.76) and pad 2 OUTPL bottom
(31.76): box 29.055-32.465 (0.5 under C17, 0.185 above C7's pad top 32.65 - place margin may want
C16 at 30.7 and C17 at 26.8). 3.3 mm from pin 24. BSPL: (30.9,27.33)-\textgreater{}(31.8,27.33)-\textgreater{}(32.4,28.9)
-\textgreater{}pad top; OUTPL: stub (30.9,27.98)-\textgreater{}(31.15,28.3)-\textgreater{}(31.15,31.3)-\textgreater{}(31.7,31.76)-\textgreater{}pad; onward east
from the pad at 31.76 to x 36.33 then the existing 36.33 vertical (shorten it to 31.76-32.65).
Rip BSPL east of 30.95, OUTPL 30.9..36.33 incl. the 36.33 vertical above 31.76.
\item C15 is the open problem: OUTNL/BSNL (OUT above BS) also needs a vertical cap or an OUT-over-top
corridor, and OUTNL's path east must pass under OUTPL's descent at x 31.15 and its 31.76 row,
i.e. at y \textgreater{}= 32.98 - into C7 (33.1) / C9 (34.6) at x 31.1-33.3 and the +12V trunk at 33.1. Two
ways: (1) move C7/C9 +12V decouplers 1 mm down or 1 mm west (they hang off the x 31.5 +12V
drop; check what sits at y 35.5) and put C15 vertical at (34.9, 33.8) with the OUTNL drop moved
from x 35.5 to 36.1 and OUTPL's vertical from 36.33 to 37.0 (BSPL's 37.0 vertical is gone);
(2) leave C15 where it is (2.4 mm from pin 20 already, it is not in the review's list) and only
move C16/C17 - then OUTPL's onward 31.76 row must stop before C15's pad 2 (32.8,29.4)... it
does not: C16's bottom pad at 31.76 sits 1.4 mm right of C15's pad column. Simplest (3): swap
roles - C15 vertical at (33.0,30.76) is the BSNL/OUTNL cap (OUTNL top 29.76, BSNL bottom 31.76:
OUTNL (30.9,29.28)-\textgreater{}(31.6,29.76)-\textgreater{}pad, BSNL (30.9,29.93)-\textgreater{}(31.4,30.9)-\textgreater{}(32.4,31.76)-\textgreater{}pad) and
C16 goes where C15 is not: BSPL/OUTPL is a BS-above-OUT pair so it may be HORIZONTAL rot 0 (BS
left) in the slot between C17 (bottom 28.555) and C15 (top 29.055)?? - no, 0.5 mm. Put C16
horizontal at (36.9, 29.5) instead (pads 35.9 BSPL / 37.9 OUTPL, box 35.195-38.605 x
28.825-30.175): OUTPL row (30.9,27.98)-\textgreater{}(33.6,27.98)-\textgreater{}(34.02,28.4) stays, continue (34.02,28.4)
-\textgreater{}(35.2,28.4)?? that passes over the BS pad - fine, BS pad is LEFT and OUTPL enters the right
pad from above at (37.9,28.4)-\textgreater{}(37.9,29.5) then down from the pad to the existing (38.4,32.65)
-\textgreater{}35.9 drop; BSPL: (30.9,27.33)-\textgreater{}(31.7,26.53)-\textgreater{}(33.5,26.53)-\textgreater{}(34.72,27.75)-\textgreater{}(35.3,27.75)-\textgreater{}
(35.9,28.8) - i.e. the existing BSPL row shortened, 0.2 clear of OUTPL's 28.4 row? 27.75+0.1 vs
28.4-0.312=28.088: 0.24 ok. Needs the OUTNL x=35.5 drop moved east of 38.6 or west: OUTNL onward
from C15's top pad (33.0,29.76) east at 29.76 hits C16's box (28.825-30.175) - so OUTNL drops at
x 34.5 (34.188-34.812, 0.38 clear of C16's box) down to 39.9 and east: check L1's courtyard
(36.8-44.2 x, \textasciitilde{}34-40.8 y) - the present 35.5 drop passes it, 34.5 will too. 6.2 mm to pin 24
is not within 4 - so (3) only fixes C17 and C15; (1) is the full fix and is the 45-60 min pass.
\end{itemize}
Next pass: resume from routing/wip-bsfix.kicad\_pcb (C17 done bar the 0.104 nudge), rip BSPL/OUTPL/ OUTNL/BSNL east of the stubs, decide (1) vs (3), then drc --refill, gates place/drc\_routed/dfm, verify\_all preview (8 known errors; the OUT stub positions (30.31, y) do not move).

\subsubsection*{12:08-12:14Z addendum (coordinator asked for a second go to 12:28Z): still reverted}
Live board unchanged (== routing/pre-bsfix.{*}), gates as recorded at 12:07Z. Checked before editing:

\begin{itemize}
\item "C17-only" is NOT gate-cleanable on its own: with C17 vertical anywhere in x 32-35 the present
BSPL 26.53 row and OUTPL 27.98/28.4 rows (to the far C16 at 38.4) run through its pad box, and
the only way past it is under y 28.755 - which is C15's pad 2 / the OUTNL 29.25 row. So C17
drags OUTPL, BSPL, and therefore C16 and C15, along. It is one coupled job, not three.
\item The vertical-column plan tightened (all clearances 0.2 checked by hand, none by DRC yet):
C17 vertical at (34.1, 26.85) rot 90 (box 33.425-34.775 x 25.145-28.555), OUTNR top / BSNR bottom;
C16 vertical at (33.6, 30.96) rot 90 with BSPL on TOP (33.6,29.96) / OUTPL bottom (33.6,31.96),
box 32.925-34.275 x 29.255-32.665; the gap 28.555-29.255 carries the BSPL row at y 28.855.
Columns going down from the pin stubs (all F.Cu): BSNL x 30.9 (0.2 wide), OUTNL x 31.55 (0.624),
OUTPL x 32.4 (0.624, from (30.9,27.98) east then down to 31.5, into the bottom pad at 31.96),
BSPL x 33.05 (from (30.9,27.33) east at 27.33 to 33.05, down to 28.855, east to 33.6, into the
top pad). OUTPL onward: bottom pad -\textgreater{} east at 31.96 -\textgreater{} x 36.33 -\textgreater{} the existing 36.33 vertical
(shortened to 31.96-32.65). OUTNL onward: down x 31.55 to y 32.9, east at 32.9 under C16
(0.235 clear of its bottom pad) to a new drop at x 35.5 (the present one, extended up from 39.9
to 32.9 - rip the 29.25 row and the 35.5 drop above 32.9).
\item The two things that block OUTNL at 32.9 and BSNL: C7 (33.1) / C9 (34.6) pads top 32.65 at
x 31.1-33.3 -\textgreater{} move C7 to (32.2,34.1) and C9 to (32.2,35.6) (check what is at y 36.0-36.5
first; L1's courtyard starts x 36.8 so it is clear) and re-hang their +12V: rip the +12V
(30.084,33.1)-(31.5,33.1)-(31.5,34.6) pieces and add (29.72,32.736)-\textgreater{}(29.72,33.7)-\textgreater{}(30.12,34.1)
-\textgreater{}(31.5,34.1)-\textgreater{}(31.5,35.6) 0.882. Then BSNL: it is fenced by OUTNL (east) and +12V (south) on
F.Cu, so it takes a via pair to B.Cu (BS nets carry mA; check\_current only polices OUT/+12V):
via at (30.9,31.9) down, under the OUTNL/OUTPL columns on B.Cu (the GND pour heals around two
0.6 vias), up at (35.0,33.6) into C15. C15 then vertical at (35.0, 34.6)?? - that sits on the
OUTNL 35.5 drop; use C15 vertical at (34.4,34.6) rot 90, OUTNL top pad (34.4,33.6) fed from the
32.9 row at x 34.4 (the row then continues to the 35.5 drop from the pad's right side? no - a
0.624 track cannot leave a 1.35-wide pad sideways without passing BSNL; so end the 32.9 row IN
the top pad and start the drop from the pad's right edge (34.9,33.6)-\textgreater{}(35.5,33.6)-\textgreater{}(35.5,39.9),
BSNL bottom pad (34.4,35.6) via at (34.4,36.5) from B.Cu). Distances: C16 3.7 mm from pin 24,
C17 4.6 mm from pin 26 (x 34.1 is 0.6 further than the 4 mm target - use x 33.7 for C17 if
BSPL's column moves to 32.95 and OUTPL's to 32.3; margins then 0.19-0.2, DRC decides).
\item Estimated at 45-60 min of place\_edit/route\_edit/DRC loops; not startable in the 17 min left at 12:11Z,
so nothing was applied. routing/wip-bsfix.kicad\_pcb still holds the first C17 attempt at x 33.0.
\end{itemize}
\subsection*{12:24-12:34Z: third C16/C17 pass - not attempted (25 min window), board unchanged}
Live board is still routing/pre-bsfix.{*} (4x4 EP, place/drc\_routed/dfm PASS). Established:

\begin{itemize}
\item C17 at x 33.0 cannot work by arithmetic: its pad box starts at x 32.325, so OUTPL's F.Cu
descent must sit at x \textless{}= 31.8 (edges 31.488-32.112). Between the pin-pad ends (30.585) and that
column only 0.9 mm remains; OUTNL's 0.624 column plus BSNL's stub (edge 31.0) need \textasciitilde{}1.4 mm.
A BSNL via at the pin is impossible (a 0.6 via inside the 0.65 pitch hits pins 19/21).
Execute the addendum's plan instead: C17 at x \textasciitilde{}34.1, C16 at x \textasciitilde{}33.6.
\item B.Cu in x 30-36, y 25-36 is bare GND pour (no tracks; one GND via at (34.3,34.9)), so the BS
nets can via there freely; that lets BSPL drop its F.Cu column and pulls C17 to about x 33.7.
\item Pocket below C16: C6 (30.6,36.2), C11 (34.15,37.3), +12V trunk at x 29.9 y 35.1-37.8 and along
y 38.089; the free slot for C7/C9 is x 32.2-33.0.
\item The full coupled pass is estimated at 45-60 min: give it a whole iteration, started first.
\end{itemize}
\subsection*{12:36-12:56Z: C16/C17 bootstrap caps moved - DONE, live board changed, gates green (hard stop 13:18Z)}
Live board = routing/post-bsfix2.kicad\_pcb (dru unchanged; fallback routing/pre-bsfix.{*} kept). Gates on it (state.json + reports/gate-\emph{.json, copies routing/bsfix2\_gate-}.json): place PASS 0/0, drc\_routed PASS 0/0 (parity + all track errors, 0 unconnected), dfm PASS 0/0. Final DRC routing/bsfix2\_g.drc.json = 0 violations, 0 warnings. Nothing committed, nothing ordered.

\subsubsection*{Where things are now (all F.Cu parts, same side as before)}
\begin{itemize}
\item C17 (33.9, 26.85) rot 90: OUTNR top pad (33.9,25.85), BSNR bottom pad (33.9,27.85).
BSNR pad to pin 26 centre 4.56 mm, to the pin-pad end (30.585,26.02) 3.79 mm (was 8.8).
\item C16 (33.7, 30.8) rot -90: BSPL top pad (33.7,29.8), OUTPL bottom pad (33.7,31.8).
BSPL pad to pin 24 centre 4.68 mm, to the pad end 3.98 mm (was 9.4).
\item C15 (32.4, 34.15) rot 0, horizontal: BSNL left pad (31.4,34.15), OUTNL right (33.4,34.15).
It had to leave (32.8,30.4) because C16 sits there now. BSNL pad to pin 20 \textasciitilde{}4.5 mm (was 2.4).
\item C7 (28.2, 33.95) and C9 (27.85, 35.4) rot 0, moved from x 32.2 to the free area under U1's
body (x 26.8-29.3, y 33.5-35.85): the pocket east of the pins is too small for C7, C9 and C15
with check.place's 0.25 mm pad-field margin (pad boxes need 0.5 mm between parts). C11 stays.
\item Refdes moved (silk over pad / silk overlap warnings): C6, C7, C9, C15, C16, C17, C18.
\item GND stitch via (34.3,34.9) removed (sat under C15's OUTNL pad); a new one at (32.6,37.1)
keeps the F.Cu GND pocket around C11's pad tied to B.Cu (without it DRC reports the island).
\end{itemize}
\subsubsection*{Routing (routing/bsfix2\_rip.json 29 rips, bsfix2\_move.json, bsfix2\_route.json + fix1..4)}
\begin{itemize}
\item OUTNR: untouched - the existing 25.42 row now runs straight through C17's top pad.
\item BSNR: (30.9,26.02)-\textgreater{}(32.2,26.02)-\textgreater{}(32.9,26.72)-\textgreater{}(32.9,27.85)-\textgreater{}pad, all F.Cu 0.2.
\item BSPL: stub -\textgreater{} via (31.35,27.2) -\textgreater{} B.Cu diagonal -\textgreater{} via (34.9,29.8) -\textgreater{} F.Cu into the top pad.
B.Cu there is bare GND pour; the pour heals around the two 0.6/0.3 vias.
\item OUTPL (0.624 F.Cu): (30.9,27.98)-\textgreater{}(31.9,27.98)-\textgreater{}(32.55,28.63) column to 31.5 -\textgreater{} (32.85,31.8)
-\textgreater{} bottom pad -\textgreater{} (35.48,31.8)-\textgreater{}(36.33,32.65) -\textgreater{} the untouched 36.33/38.4 path to L1.
\item OUTNL (0.624 F.Cu): (30.9,29.28)-\textgreater{}(31.7,29.28) column to 33.0 -\textgreater{} row to (35.5,33.0) -\textgreater{} the
35.5 drop to 39.9 (re-added from 33.0). Stub (33.4,33.0)-\textgreater{}(33.4,34.15) into C15's pad.
\item BSNL: all F.Cu now, no via: (30.9,29.93)-\textgreater{}(30.85,30.3)-\textgreater{}(30.85,34.15)-\textgreater{}(31.4,34.15), running
0.165 mm right of the +12V pin-pad ends (30.585). A via there is impossible: U1's pads are
1.731 mm long along x (their ends are at 30.585, not 29.89).
\item +12V: C7/C9 hang re-done: (29.72,32.736)-\textgreater{}(29.9,32.916)-\textgreater{}(29.9,35.109) feeds C6 as before;
(27.5,32.736) drop off the west trunk to C7 pad 1, then (27.15,34.3)-\textgreater{}(27.15,35.4) to C9.
\item Clearance is 0.1524 (netclass default 0.152), not the 0.2 the earlier passes assumed - that
is what made the two columns fit between the pad ends and C16.
\end{itemize}
\subsubsection*{Left}
\begin{itemize}
\item verify gate still FAIL with the 8 known errors (not re-run; OUT stub positions unchanged).
\item check\_current's "6 vias per transition" only polices OUT/+12V; the BS nets carry mA.
\item The next order needs a fresh fab export (board changed after reports/fab\_export.json).
\end{itemize}
\section{How the Design Process Went}
Each phase the run passed through, top to bottom; a red box is a point where it went back to an earlier phase, with the first decision that sent it there. The current phase has the red outline.
\begin{center}
\includegraphics[width=0.8\textwidth]{reports/design_full/flow.pdf}\\
{\small\texttt{reports/design\_full/flow.pdf}}
\end{center}
\section{Verification}
\emph{Pending \textemdash{} produced at P8.}
\section{DFM and Fabrication}
\emph{Pending \textemdash{} produced at P9.}
\section{Run Record (Appendix)}
\subsection*{Phase digests}
\paragraph*{P0}
\begin{quote}\small\ttfamily\raggedright
\# P0 digest\\
- requirements.md written by requirements-analyst (fable), check\_requirements pass, no mode token.\\
- 12 open questions answered by the orchestrator as delegate (defaults accepted); Q6 current 3 A PROVISIONAL.\\
- Size soft target \textless{}= 44 x 34 mm, 2 layers; thermal pad / heat-sinking outranks size.\\
- Assembly ASSUMED JLC PCBA top side; qty 5; bare wire pads J1-J4 + TP1; no protection, no holes.
\end{quote}
\paragraph*{P1}
\begin{quote}\small\ttfamily\raggedright
\# P1 digest\\
- Roster: component-scout (sonnet), reference-design (sonnet, incl. both reference-design sources), power-architect (opus).\\
- DEVIATION: TPA3116D2 is pad-UP only (DAD); swapped to pin-compatible pad-down TPA3118D2DAPR (C46497) so the board copper sinks it.\\
- +12V 4 A (3.6 A at 10\% THD, 4 ohm); outputs 3 A; chip loss \textasciitilde{}3 W (4 ohm 1\%), \textasciitilde{}1.2 W (8 ohm), \textasciitilde{}0.8 W music.\\
- 44x34 mm holds 8 ohm full power and 4 ohm music; sustained 4 ohm full sine needs \textasciitilde{}50x50 mm (not required).\\
- L = SLO0630H100MTT 7.3x6.6 mm (Isat 5.5 A) - the size driver; 680 nF filter caps; 3.3R+10nF damping; 2x100 uF (polymer+electrolytic).\\
- 2 oz copper preferred; PLIMIT-\textgreater{}GVDD, FAULTZ-\textgreater{}SDZ, 20 dB gain (5.6k), 100k pull-ups.\\
- Quilter source had no board-specific content; JLC source: 44x34, 2L, 45 fp / 40 fitted, 0 ERC/DRC, never fabricated.
\end{quote}
\paragraph*{P2}
\begin{quote}\small\ttfamily\raggedright
\# P2 digest\\
- Architect (fable/high): 4 blocks - B1 12 V/4 A wire-pad input + 100 uF polymer \& 100 uF electrolytic bulk; B2 stereo line-in, 2.2 uF X7R coupling; B3 TPA3118D2DAPR BTL, BD, 400 kHz, 20 dB, master; B4 10 uH + 680 nF LC per leg, 3.3R+10nF damping.\\
- One schematic sheet, no hierarchy (38 fitted parts); canonical net names in sheets.md, used by constraints.json.\\
- Stackup: 2 layers, 2 oz outer preferred; B.Cu one GND pour with the PowerPAD thermal-via array into it.\\
- Size driver: four 7.3x6.6 mm inductors; target \textless{}= 44x34 mm, placement earns the outline (provisional, then fit).\\
- constraints\_lint exit 0 (11 unitless operating-point warnings, informational).\\
- Coverage P2: 4/4 slots gap (no checklists for these topologies) -\textgreater{} 4 research tasks opened automatically.
\end{quote}
\paragraph*{P3}
\emph{(no digest recorded)}
\paragraph*{P4}
\emph{(no digest recorded)}
\paragraph*{P5}
\emph{(no digest recorded)}
\paragraph*{P6}
\emph{(no digest recorded)}
\paragraph*{P7}
\emph{(no digest recorded)}
\subsection*{Gate attempt history}
{\small
\begin{longtable}{llll}
\hstoprule
\hshead{Timestamp} & \hshead{Gate} & \hshead{Status} & \hshead{Failing/Total} \\
\hline
\endhead
2026-09-27T03:19:32 & erc & pass & 0/0 \\
2026-09-27T03:45:59 & erc & pass & 0/0 \\
2026-09-27T04:10:46 & place & fail & 49/49 \\
2026-09-27T04:11:28 & place & fail & 49/49 \\
2026-09-27T04:11:48 & place & fail & 49/49 \\
2026-09-27T04:13:04 & place & fail & 21/21 \\
2026-09-27T04:14:33 & place & fail & 3/3 \\
2026-09-27T04:15:31 & place & fail & 1/1 \\
2026-09-27T04:15:55 & place & fail & 1/1 \\
2026-09-27T04:16:16 & place & pass & 0/0 \\
2026-09-27T04:17:57 & place & pass & 0/0 \\
2026-09-27T04:19:59 & place & pass & 0/0 \\
2026-09-27T04:52:12 & drc\_routed & fail & 56/56 \\
2026-09-27T05:29:04 & place & fail & 5/5 \\
2026-09-27T05:47:45 & place & fail & 2/2 \\
2026-09-27T05:51:46 & place & fail & 1/1 \\
2026-09-27T05:52:46 & place & fail & 1/1 \\
2026-09-27T05:53:10 & place & pass & 0/0 \\
2026-09-27T05:54:01 & place & pass & 0/0 \\
2026-09-27T05:54:06 & drc\_routed & fail & 7/7 \\
2026-09-27T05:54:43 & place & pass & 0/0 \\
2026-09-27T05:54:48 & drc\_routed & fail & 6/6 \\
2026-09-27T05:55:08 & place & pass & 0/0 \\
2026-09-27T06:01:34 & place & fail & 5/5 \\
2026-09-27T06:02:16 & place & pass & 0/0 \\
2026-09-27T06:02:17 & place & pass & 0/0 \\
2026-09-27T06:02:23 & drc\_routed & fail & 6/6 \\
2026-09-27T06:23:13 & place & pass & 0/0 \\
2026-09-27T06:23:17 & drc\_routed & fail & 3/3 \\
2026-09-27T06:24:44 & drc\_routed & fail & 2/2 \\
2026-09-27T06:25:19 & drc\_routed & pass & 0/0 \\
2026-09-27T06:25:22 & place & pass & 0/0 \\
2026-09-27T06:34:02 & place & pass & 0/0 \\
2026-09-27T06:34:07 & drc\_routed & pass & 0/0 \\
2026-09-27T06:35:39 & dfm & pass & 0/0 \\
2026-09-27T06:37:18 & place & pass & 0/0 \\
2026-09-27T06:37:23 & drc\_routed & pass & 0/0 \\
2026-09-27T06:38:29 & dfm & pass & 0/0 \\
2026-09-27T06:55:35 & verify & fail & 8/31 \\
2026-09-27T06:57:55 & verify & fail & 8/31 \\
2026-09-27T07:21:41 & drc\_routed & pass & 0/0 \\
2026-09-27T07:21:43 & place & pass & 0/0 \\
2026-09-27T07:26:10 & dfm & pass & 0/0 \\
2026-09-27T07:26:25 & verify & fail & 8/31 \\
2026-09-27T07:49:27 & place & pass & 0/0 \\
2026-09-27T07:50:41 & drc\_routed & pass & 0/0 \\
2026-09-27T07:54:33 & dfm & pass & 0/0 \\
2026-09-27T08:07:10 & place & pass & 0/0 \\
2026-09-27T08:07:14 & drc\_routed & pass & 0/0 \\
2026-09-27T08:13:09 & dfm & pass & 0/0 \\
2026-09-27T08:13:19 & verify & fail & 8/31 \\
2026-09-27T08:51:08 & place & fail & 7/7 \\
2026-09-27T08:51:11 & drc\_routed & pass & 0/0 \\
2026-09-27T08:51:21 & dfm & pass & 0/0 \\
2026-09-27T08:54:10 & place & pass & 0/0 \\
2026-09-27T08:55:02 & place & pass & 0/0 \\
2026-09-27T08:55:04 & drc\_routed & pass & 0/0 \\
2026-09-27T08:55:11 & dfm & pass & 0/0 \\
2026-09-27T08:57:46 & dfm & pass & 0/0 \\
2026-09-27T08:57:51 & verify & fail & 8/32 \\
\hline
\end{longtable}
}
\subsection*{Issues}
\emph{(no entries)}
\subsection*{Human checkpoints}
{\small
\begin{longtable}{lllp{2.6cm}p{6.4cm}}
\hstoprule
\hshead{Id} & \hshead{Phase} & \hshead{Status} & \hshead{When} & \hshead{Note} \\
\hline
\endhead
1 & P2 & approved & 2026-09-27T03:07:31 & delegate (owner: decide for yourself). Coverage: C46497 covered; B2/B3/B4 provisional, all 4 research tasks closed verified; B1 gap accepted by decision. TPA3118D2 DAP pad-down instead of 3116 DAD pad-up (needs heatsink); 2L 2oz; 4 A input; riskiest: 2L thermal at 4 ohm sustained, hot-plug ring margin 1.15x. \\
2 & P4 & approved & 2026-09-27T03:46:57 & delegate (owner: decide for yourself). Schematic 44 parts/29 nets, ERC 0/0, netlist\_audit pass; review 0 err/5 warn, W1-W3 fixed, W4 silk legend at layout, W5 X5R accepted. Deviation from the reference design: TPA3118D2 pad-down + SMBJ15A TVS added + 2x100uF bulk. \\
3 & P6 & approved & 2026-09-27T04:35:06 & delegate (owner: decide for yourself). Place gate PASS, DRC 0 after silk fix, 53.6x36.8 mm 2L accepted as deviation (decision). \\
4 & P8 & not reached &  &  \\
5 & P10 & not reached &  &  \\
\hline
\end{longtable}
}
\subsection*{Generation cost}
Model cost of generating this board: \textbf{\$97.57}.

{\small
\begin{longtable}{p{10cm}r}
\hstoprule
\hshead{Step} & \hshead{Cost} \\
\hline
\endhead
Before the workspace (brief and skill read) & \$0.50 \\
Brief and intake & \$1.16 \\
Research & \$3.84 \\
Architecture & \$14.22 \\
Schematic & \$3.10 \\
Placement & \$4.88 \\
Routing & \$0.42 \\
Verification & \$0.51 \\
DFM, order files and guide & \$0.11 \\
Not in the transcript (side calls) & \$68.83 \\
\hline
\end{longtable}
}
\section{Artifact Index}
{\small
\begin{longtable}{p{9cm}rp{4cm}}
\hstoprule
\hshead{File} & \hshead{Size} & \hshead{Note} \\
\hline
\endhead
\texttt{state.json} & 69,768 B &  \\
\texttt{kicad/PCB-0017-A\_stereo-class-d-amp.kicad\_pcb} & 261,364 B &  \\
\texttt{kicad/PCB-0017-A\_stereo-class-d-amp.kicad\_sch} & 151,584 B &  \\
\texttt{fab/PCB-0017-A\_stereo-class-d-amp\_gerbers.zip} & 45,371 B &  \\
\texttt{fab/BOM.csv} & 743 B &  \\
\texttt{fab/CPL.csv} & 1,334 B &  \\
\texttt{brief/brief.md} & 953 B &  \\
\texttt{requirements.md} & 9,052 B &  \\
\texttt{architecture/blocks.md} & 8,474 B &  \\
\texttt{architecture/constraints.json} & 15,472 B &  \\
\texttt{architecture/power\_tree.md} & 6,460 B &  \\
\texttt{architecture/sheets.md} & 6,644 B &  \\
\texttt{architecture/stackup.md} & 4,159 B &  \\
\texttt{reports/PCB-0017-A\_stereo-class-d-amp\_bottom.png} & 38,076 B &  \\
\texttt{reports/PCB-0017-A\_stereo-class-d-amp\_top.png} & 230,278 B &  \\
\texttt{reports/board\_init.json} & 1,038 B &  \\
\texttt{reports/bom\_cpl.json} & 22,913 B &  \\
\texttt{reports/check\_thermal.json} & 2,013 B &  \\
\texttt{reports/cost.json} & 6,761 B &  \\
\texttt{reports/fab\_export.json} & 2,359 B &  \\
\texttt{reports/gate-dfm.json} & 1,552 B &  \\
\texttt{reports/gate-drc\_routed.json} & 1,022 B &  \\
\texttt{reports/gate-erc.json} & 822 B &  \\
\texttt{reports/gate-place-optionA.json} & 4,103 B &  \\
\texttt{reports/gate-place-optionA3.json} & 4,171 B &  \\
\texttt{reports/gate-place.json} & 1,407 B &  \\
\texttt{reports/gate-verify.json} & 8,573 B &  \\
\texttt{reports/guide\_facts.json} & 12,627 B &  \\
\texttt{reports/placement-notes.md} & 2,958 B &  \\
\texttt{reports/review-schematic.json} & 6,263 B &  \\
\texttt{reports/review-schematic.md} & 14,056 B &  \\
\texttt{reports/review-verify.md} & 9,422 B &  \\
\texttt{reports/schematic.pdf} & 149,082 B &  \\
\texttt{reports/top.net} & 74,740 B &  \\
\texttt{reports/verify-waivers.proposed.json} & 6,139 B &  \\
\hline
\end{longtable}
}
\end{document}
